\begin{center}
{\Large\bfseries The denominator estimate at the two ramified primes}\\[4pt]
Edition 0.81 \quad 25 November 2026
\end{center}

For a good point at either distinguished prime $q,d$, we bound the
congruence denominator by the Selmer, character and class-number
terms in the original family inequality. The character bound comes
from an integral estimate on the cocycle of a crystalline extension.
The rank estimate uses the two Hodge orientations over each residue
embedding. The opposite orientations at $d$ require an additional
integral coefficient calculation. With the global determinant
argument these estimates prove the thirteenth of the fifteen fields
of the denominator family conclusion. The final sections carry the
character estimate to the middle range and prove integral lifting at
the split primes there.

\paragraph*{The bound at the two primes.}
Let $u=\log b$ and $B=b/p+S_c+t+\eta+u$.
On the original family data, for each sufficiently late qualifying member,
twisted index $i$ at $p=q,d$ and good point $\pi$, we obtain
\[
\begin{aligned}
 v_{\rm cc}(i,\pi)\le {}&
 s_\pi+\sum_{1\le j<t} u_j+v_p(h_E)+p^{-2}\\
 &{}+2080 f(rB+mb)+181 f^2t^2(r+t)u .
\end{aligned}
\]
The first line records the Selmer length, the positive character twists,
the class-number term and the auxiliary correction. To control the
remaining rank contribution, restrict the global rank classes at all
places over $p$ and take their images modulo the local finite classes.
The kernel contributes to $s_\pi$. The image has length bounded by
the sum of the local lifting estimates. In particular the estimate
must hold for all the integral maps tested by the global count.
The calculation giving the displayed constants is J81H1--J81H2.
The last two family fields, for the middle range and the small-prime
correction, are still open.

\begingroup\emergencystretch=3em
\section{The integral torsion exponent for the truncated period module}

% BEGIN W81T1
\paragraph{W81T1. The prescribed integral quotient and its single graded pieces.}
Write $p$ for the local prime, $K_j=K_0(\mu_{p^j})$,
$G_j=G_{K_j}$, $H=G_{K_\infty}$ and $C=\mathbb C_p$. Use the
chosen tied Fontaine and Wach records of L78W13 on a good layer
$(K,s,L)$. Thus $K\subset K_s$, $s\ge1$ and $s\ge2$ if $p=2$.
Let $T$ be the underlying finite free $\mathbb Z_p$-lattice of
a source object satisfying $\mathrm{HypW}(M_0,h_0)$.
The integral coefficient module on a window $a<b$ is
\[
 X_{a,b}=T\otimes_{\mathbb Z_p}
         (\xi^a A_{\inf}/\xi^b A_{\inf}).
\]
Integral negative powers here mean the corresponding fractional
modules inside the large Witt ring, where $\xi$ is a unit.
Multiplication by $\xi^{-a}$ identifies their underlying quotients
with $A_{\inf}/\xi^{b-a}$. This is the actual $Q_{a,b}$ with its
transported Galois action, including the action on $\xi^a$ itself.
LS4--LS5 give $p$-torsion-freeness, $p$-adic completeness, and
after inverting $p$ the identification with the corresponding
filtered quotient $\mathrm{Fil}^aB_{\rm dR}/\mathrm{Fil}^bB_{\rm dR}$.
Indeed completion of $A_{\inf}[1/p]$ along the regular element
$\xi$ gives precisely its original quotients by $\xi^j$ for each
finite $j$. Tensoring with the finite free $T$ preserves these
statements and the exact sequences cutting a window at any integer.

For $g\in G_s$, $\theta(g\xi/\xi)=\chi(g)$. In fact
$t=\log[\epsilon]$ satisfies $gt=\chi(g)t$. Since
$[\epsilon]-1=\xi([\epsilon_1]-1)$, the residue of $t/\xi$ is
$\zeta_p-1$, which is nonzero and fixed by $g$. Reduce the identity
$(g\xi/\xi)g(t/\xi)=\chi(g)t/\xi$ in the de Rham DVR.
The ratio $g\xi/\xi$ is an integral unit since $g$ preserves the
kernel of $\theta$. Division by $\xi^i$ followed by the
surjective $\theta:A_{\inf}\to O_C$ therefore identifies
$X_{i,i+1}$ with $T\otimes O_C(i)$, integrally and equivariantly.
This also explains the positive sign of $i$ in this coefficient twist.
% END W81T1

% BEGIN W81T2
\paragraph{W81T2. Integral approximation to a primitive on one graded piece.}
Let $M(\sigma)$ be the action on the chosen $\mathbb Z_p$-basis of
$T$ of rank $d$, and put
\[
 m=2s+1+\min\{n\ge0:M_0+h_0\le p^n(p-1)\},\quad
 J=\mathrm{Jset}(T),\quad
 B=O_C^d.
\]
We use all the components of the source filtration. In particular
$J$ contains zero whenever zero is an actual jump. A79L6 constructs
on this very $\mathrm{HypW}$ lattice an invertible matrix $D$ with
$pB\subseteq DB\subseteq B$ and
$U(\sigma)=D^{-1}M(\sigma)\sigma(D)\in\mathrm{GL}_d(O_{K_m})$
on $G_m$. The representation $U$ is trivial on $H$, and there is
a $g\in G_m$ with $v_p(\chi(g)-1)=m$ and
$\|U(g)-1\|<|\zeta_p-1|$. A79L7, before its projection to one
embedding, gives a $C$-basis of vectors $b$ of weights $u\in J$
satisfying $\chi(\sigma)^uM(\sigma)\sigma(b)=b$ on $G_m$.
Here the required de Rham comparison on the given good layer
follows from its restricted crystallinity as in A79L7.
The weights lie in the finite interval $[-M_0,h_0]$.

Fix an integer $i$ and act on $C^d$ by
$\sigma*x=\chi(\sigma)^i M(\sigma)\sigma(x)$.
Suppose $x\in C^d$ has $\sigma*x-x\in B$ on $G_m$.
In $U$-coordinates let $y=pD^{-1}x$. Its differences for the
action with $U$ are integral since $pD^{-1}B\subset B$.
On $H$ this action is coordinatewise, so A79L2 gives
$py=z+b_0$ with $z\in K_\infty^d$, $b_0\in B$.
Let $A=\chi(g)^i U(g)$ and $\Gamma=g|_{\widehat K_\infty}$.
Then $(A\Gamma-1)z$ is integral. Split
$z=z_f+z_c$ in the normalized-trace splitting of A79L3, into
$K_m^d$ and the rational complementary summand. Both projections
are integral and commute with $A\Gamma$. In particular
$(A-1)z_f$ is integral and $p z_c$ is integral by the complementary
inverse of A79L3, since $p/(\zeta_p-1)$ is integral.

The vectors $D^{-1}b$ of weight $u$ for $U$ lie in
$\widehat K_\infty^d$ by $H$-invariance and A79L2. Their
complementary parts vanish by the injectivity of
$\chi(g)^uU(g)\Gamma-1$ there, from A79L3. They are thus eigenvectors
over $K_m$ for $A$ with eigenvalues $\chi(g)^{i-u}$; in particular
they span $C^d$. This argument retains the eigenspace with $u=i$.
% END W81T2

% BEGIN W81T3
\paragraph{W81T3. The resonant weight costs no denominator.}
Put
\[
 F_i(Z)=\prod_{u\in J,\ u\ne i}(Z-\chi(g)^{i-u}),\qquad
 r_i=\sum_{u\in J,\ u\ne i}(m+v_p(u-i)),\quad n_i=\max(1,r_i).
\]
The polynomial has integral coefficients and
$(A-1)F_i(A)=0$ by the basis in W81T2. The binomial valuation
calculation in A79L4 gives $v_p(F_i(1))=r_i$, with $F_i(1)=1$
when the product is empty. As $F_i(Z)-F_i(1)$ is an integral
multiple of $Z-1$, the integrality of $(A-1)z_f$ shows that
\[
 z_* = \frac{p^{n_i}}{F_i(1)}F_i(A)z_f
       \quad\text{satisfies}\quad p^{n_i}z_f-z_*\in O_{K_m}^d .
\]
We have $Az_*=z_*$. It is fixed for the twisted $U$-action by $H$
and by $g$, because its coordinates lie in $K_m$. The image of
$g$ generates a dense subgroup in $\mathrm{Gal}(K_\infty/K_m)$:
$v_p(\chi(g)-1)=m$, with $m\ge2$ at $p=2$, makes its successive
$p$-powers generate each successive quotient. Continuity therefore
gives invariance under all of $G_m$.

Now $p^{n_i}z_c\in B$, and $py=z+b_0$ in W81T2. Multiplication
back by $D$ shows
\[
                  p^{2+n_i}x-Dz_*\in B .
\]
The correction $Dz_*$ is invariant for the twisted original action.
By W81T1 this proves the following assertion about the actual
graded piece: every rational primitive with integral differences
becomes integral modulo a rational invariant after multiplication
by $p^{2+n_i}$. In particular a torsion continuous cocycle of
this piece is killed in cohomology by that power: divide the
integral primitive of a killing multiple by this multiple and use
the assertion just proved. The use of distinct weights in $F_i$
makes the exponent independent of their multiplicities.
% END W81T3

% BEGIN W81T4
\paragraph{W81T4. Lifting rational invariants without increasing the exponent.}
For $a\le j\le b$ the map $X_{a,b}[1/p]\to X_{a,j}[1/p]$
is surjective on $G_m$-invariants. To see this, write
$V=T[1/p]$ and take the adapted basis of invariant de Rham periods
$w_k$ from A79L7, of filtration degrees $u_k$. Its leading vectors
$b_k=t^{-u_k}w_k\bmod t$ form a $C$-basis as in W81T2.
The invariants in $V\otimes\mathrm{gr}^q B_{\rm dR}$
are precisely the $K_m$-span of the leading symbols of those
$w_k$ with $u_k=q$. Indeed after extracting the formal factor $t^q$
expand in the basis $b_k$. The coefficient at $k$ then satisfies
$\sigma(c)=\chi(\sigma)^{u_k-q}c$ on $G_m$. It is fixed by $H$
and hence belongs to $\widehat K_\infty$ by A79L2. The complementary
part in A79L3 is zero by its scalar twisted inverse. On $K_m$ the
generator $g$ forces $c=0$ unless $u_k=q$, because $\chi(g)$ has
infinite order. For equal weights the coefficient can be any
element of $K_m$, embedded in the de Rham period ring in the
given way.

Start with an invariant class modulo $\mathrm{Fil}^j$ at filtration
$a$. Its image in degree $a$ is invariant. Subtract the just-described
linear combination of actual invariant periods of degree $a$ to
increase the filtration by one. Repeat through degree $j-1$.
The sum of the subtracted periods gives an invariant representative
in $V\otimes\mathrm{Fil}^a B_{\rm dR}$. Its reduction modulo
$\mathrm{Fil}^b$ gives the desired lift by W81T1.
In particular an integral invariant class lifts after multiplication
by some power of $p$ to an integral invariant class. This power is
used only to identify rational invariant subspaces; it is never
added to the numerical torsion exponent.
% END W81T4

% BEGIN W81T5
\paragraph{W81T5. Adding the graded exponents on the full window.}
For a stable $p$-torsion-free module $X$ call a rational vector
$x\in X[1/p]$ integral modulo invariants at exponent $c$ if
$p^cx-x_0\in X$ for some $x_0\in X[1/p]^{G_m}$.
In an equivariant exact sequence
$0\to X'\to X\xrightarrow{q}X''\to0$ of the present modules,
suppose the assertion for every vector with integral differences
holds at exponents $c'$ and $c''$ on the two ends, and $q$ is
surjective on rational invariants. For such a vector $x$ in the
middle, we may write $p^{c''}q(x)=u''+z''$ with $u''$ invariant
and $z''$ integral. Lift $u''$ to a rational invariant $u$ and
$z''$ to an integral element $z$. Then
$y=p^{c''}x-u-z\in X'[1/p]$ has integral differences in $X'$:
they belong to $X\cap X'[1/p]=X'$, since $X''$ is torsion-free.
Apply the assertion at $c'$ to $y$ and add $p^{c'}z$ back. This
proves the middle assertion with exponent $c'+c''$.
The proof does not choose an integral lift of an invariant with a
specified bound.

Apply this argument to
$0\to X_{b-1,b}\to X_{a,b}\to X_{a,b-1}\to0$.
W81T4 gives precisely its rational invariant lifting and W81T3
is the one-step case. Thus every rational primitive with integral
differences on $G_m$ is integral modulo invariants at exponent
$C_{a,b}=\sum_{a\le i<b}(2+n_i)$.
To descend to $G_s$, use a set $R$ of representatives of the
left cosets $G_s/G_m$. For a rational primitive $x$ with integral
differences on $G_s$, choose $u\in X_{a,b}[1/p]^{G_m}$ with
$p^{C_{a,b}}x-u$ integral. The sum $\sum_{r\in R}r u$ is
$G_s$-invariant, and
$[G_s:G_m]p^{C_{a,b}}x-\sum_r r u$ is integral since each
$p^{C_{a,b}}(x-rx)$ is integral. Eisenstein's cyclotomic degree
formula over the unramified $K_0$ gives
$[G_s:G_m]=[K_m:K_s]=p^{m-s}$, including $m=s$.
The difference of the $p$-valuations of their absolute degrees
is therefore $m-s$, exactly $\mathrm{vdegMs}(m)$.
% END W81T5

% BEGIN W81T6
\paragraph{W81T6. The torsion multiplier consumed by the ramified extension.}
Fix any continuous cocycle $c$ of $G_s$ in $X_{a,b}$ for which
$p^kc$ is an integral coboundary. If $v$ is its primitive, then
$x=p^{-k}v$ has differences $c$ because $X_{a,b}$ is
$p$-torsion-free. W81T5 now gives an integral primitive of
$p^{N_{a,b}}c$ with
\[
 N_{a,b}=m-s+
       \sum_{a\le i<b}\left(2+\max\left(1,
               \sum_{\substack{u\in J\\u\ne i}}
                      (m+v_p(u-i))\right)\right).
\]
This is the source's $\mathrm{cLL}$ for all $a<b$ on the chosen
records. In particular in the extension-class argument take $a=0$,
$b=e+1$, $T$ the kernel lattice, and
$c(g)=f(g)\otimes[1]$ on $G_s$. Its assumed torsion and its
cocycle identity suffice to obtain the actual integral coboundary
in $T\otimes Q_{0,e+1}$ after multiplication by $p^{N_{0,e+1}}$.
The additional terms in $\mathrm{cTB}$ belong to the subsequent
passage from this coboundary to Bloch--Kato finiteness.
% END W81T6

\section{From the integral connecting cocycle to the ramified character condition}

% BEGIN K81C1
\paragraph{K81C1. A torsion class on the truncated integral period ring.}
Let $p$ be odd, $K_0$ the maximal unramified subfield of $K$, and
$F=K_0(\mu_p)\supset K$. Use the level-one data on the original tied
period rings. Denote the cyclotomic character by $\chi$, the subgroup
$G_F$ by $G_1$, and $G_{K_0(\mu_{p^\infty})}$ by $H$. Fix arbitrary
$e\ge0$ and put $M=e+1$, $A=A_{\inf}$, $A_M=A/(\xi^M)$ and
$B_M=B_{\rm dR}^+/\mathrm{Fil}^{M}$. Then $A_M[1/p]=B_M$ on
these rings: completion in the principal ideal $(\xi)$ induces an
isomorphism on each finite quotient, and $\xi$ generates the kernel
of $\theta$ after inverting $p$. Before inverting $p$ the quotient
embeds, since $p$ is regular modulo $\xi$ and hence modulo its
powers by induction. The map carries $1$ to $1$ and is equivariant.

Take a tested integral short exact sequence $Y\to\widetilde Y\to Z$
from $\mathrm{TBat}$ with trivial quotient line and let $y$ be the
given integral lift. Write $f(g)$ for the unique lattice vector
mapping to $gy-y$. Injectivity gives $f(gh)=f(g)+gf(h)$ and a
linear left inverse to the embedding shows continuity. Over $F$ the
middle term is crystalline by hypothesis. Strictness at filtration
zero (J80H1, or the strict filtered de Rham comparison) gives a
$G_1$-fixed vector in $\widetilde Y\otimes B_{\rm dR}^+$ above
$\bar y\otimes1$. Subtract it from $y\otimes1$. The result lies in
$Y\otimes B_{\rm dR}^+$ by exactness of the underlying vector
spaces, and has coboundary $f\otimes1$. After reduction to $B_M$,
clear the finitely many $p$-denominators of its coordinates in a
$\mathbb Z_p$-basis of the kernel lattice $T$. Thus $f\otimes1$ is
a torsion cocycle in $T\otimes A_M$. Its reductions modulo open
lattice neighborhoods are continuous by continuity of $f$.

Use the integral Wach torsion bound W81T6 at $(a,b)=(0,M)$ for the
$\mathrm{HypW}$ of this actual kernel. In the notation of the
original $\mathrm{cLL}$ put
\[
 N=\mathrm{cLL}\bigl(p,m,\mathrm{vdegMs}(m),J(Y),0,e+1\bigr),
 \qquad m=2+ n_0(p,M_0+h_0)+1.
\]
It yields one $X\in T\otimes A_M$ with
$(g-1)X=p^N f(g)\otimes1$ for every $g\in G_1$. The subsequent
argument uses this identity at the stated $N$.
% END K81C1

% BEGIN K81C2
\paragraph{K81C2. Integral contraction and rational lifting.}
Let $\phi:T\to V/T_e$ be the tested equivariant map, where
$V$ is the coefficient line with character $\chi^{-e}$ and $T_e$
its prescribed lattice. Lift $\phi$ linearly to $\Phi:T\to V$
on a free integral basis. Relative to a $\mathbb Z_p$-basis of
$T_e$ every coordinate $\psi$ of $\Phi$ is $\mathbb Z_p$-linear,
bounded on $T$ and equivariant modulo $\mathbb Z_p$. Contract the
finite tensor $X=\sum t_i\otimes x_i$ with $\psi$. This gives
$b=\sum\psi(t_i)x_i\in B_M$. For $c(g)=\psi(f(g))$ and
$D_g b=\chi(g)^{-e}g b-b$, equivariance modulo $\mathbb Z_p$
gives directly
\[
                 p^N c(g)-D_g b\in A_M\qquad(g\in G_1).
\]
Indeed the error is $\sum(\psi(gt_i)-\chi(g)^{-e}\psi(t_i))gx_i$;
the arbitrary common denominator of $\psi$ costs no power of $p$.
Moreover $c$ modulo $\mathbb Z_p$ is an actual continuous cocycle.

Put $h=1+v_p(e+1)$. By integral Tate duality
$H^2(F,\mathbb Z_p(-e))$ is dual to the invariants of
$\mathbb Q_p/\mathbb Z_p(e+1)$. An element $\gamma\in G_1$
with $\chi(\gamma)=1+p$ acts on these by $(1+p)^{e+1}$, whose
difference from one has valuation $h$. Such a $\gamma$ exists
because $K_0$ is unramified and the cyclotomic tower has its full
procyclic group above level one. Hence $p^h$ kills the obstruction
to lifting the class of $c$. The exact sequence for
$\mathbb Z_p(-e)\subset\mathbb Q_p(-e)$ gives a continuous rational
cocycle $z$ with $z(g)-p^h c(g)\in\mathbb Z_p$ on $G_1$:
a discrepancy by a torsion coboundary is removed by subtracting
the coboundary of any rational lift of its value. The rational
cocycle $Z=p^N z$ consequently satisfies
$Z(g)-D_g(p^hb)\in A_M$ on $G_1$.
% END K81C2

% BEGIN K81C3
\paragraph{K81C3. The coefficient of the logarithmic cocycle and its bound.}
Write $C=\mathbb C_p$, $\delta=\zeta_p-1$, and $t=\log[\epsilon]$.
For a rational $\chi^{-e}$-cocycle $Z$ there is a unique
$\alpha\in F$ such that, in $B_M$, for some $b_0$,
\[
 Z(g)=\log\chi(g)\,\alpha t^e+D_g b_0 .
\]
We give the graded calculation and its integral estimate. For a
continuous unit-ball-valued cocycle on $H$ acting on $C$, multiply
by $p$ to obtain the coboundary of an element of the unit ball.
Indeed in a finite Galois quotient with coset representatives $s$
pick integral $a$ of trace $p$ by A79L1. The negative weighted sum
$-\sum_s s(a)c(s)$ has norm at most the bound on $c$ and its
coboundary differs from $pc$ by at most the bound on the small
open kernel, by permuting the cosets in the cocycle identity.
Continuity makes the latter bound arbitrarily small. Apply this
again to the residual cocycle divided by $p$, choosing successive
errors at most $|p|2^{-j}$; the successive increments have norm
at most $2^{1-j}$ and converge to a primitive of $pc$.

For a unit-ball cocycle of $C(j)$ on $G_1$ the twist is trivial
on $H$. Subtract this integral coboundary from $p$ times the
cocycle. The result factors through $G_1/H$, and its value on
$\gamma$ is in the unit ball of $C^H=\widehat K_\infty$.
Use A79L3's integral splitting into $F$ and the complementary
summand. On the complement $\chi(\gamma)^j\gamma-1$ has inverse
of norm at most $|\delta|^{-1}$. On $F$, if $j\ne0$, divide by
$(1+p)^j-1$, of valuation $1+v_p(j)$. This gives a primitive of
the original cocycle of norm at most $p^{2+v_p(j)}$; equality of
cocycles on $H$ and $\gamma$ implies equality by density.
There are no invariant vectors for $j\ne0$, by the same splitting.
If $j=0$ the residual $F$-coordinate is instead $p\alpha\log(1+p)$.
Thus the unique logarithmic coefficient of a bounded cocycle has
norm at most $p^2$. Uniqueness follows on projecting a putative
logarithmic coboundary to $F$: its primitive is fixed by $H$ and
normalized trace kills $\gamma-1$ there. Compactness permits
scaling any continuous cocycle to the unit ball.

Filter $B_M$ by powers of $t$. At step $i$ the degree-$i$ error
is a $C(i-e)$-cocycle. At $i<e$ lift a primitive via $\theta$
and subtract its coboundary times $t^i$; at $i=e$ the preceding
weight-zero calculation gives the displayed formula. The quotient
topology is the $p$-adic lattice topology on $A_M[1/p]$, and
its graded maps to $C$ are continuous. The Galois action on
these quotients is continuous: on $A$ this follows on finitely
many Witt coordinates modulo $(p^k,\xi^M)$ from continuity on
the tilt in its valuation topology, and extends by inverting $p$.
The same graded descent for the difference of two such expressions
proves uniqueness: the earlier degrees have zero invariants, and
a logarithmic coboundary in the last degree has coefficient zero.

Suppose also that $Z-D_g b$ is integral in $A_M$. The difference
of the two expressions says $\log\chi(g)\alpha t^e$ is an integral
error plus a coboundary with primitive in degree at least zero.
At degree $i<e$ write the primitive as $t^i y$. In the coordinate
relative to $\xi^i$ the residue is $\delta^i\theta(y)$, since
$t/\xi$ has residue $\delta$ by $\mu=\xi([\epsilon^{1/p}]-1)$
and $t/\mu\equiv1\pmod{\mathrm{Fil}^1}$. The integral error in
this coordinate is in $\mathcal O_C$ because
$A\cap\mathrm{Fil}^i=\xi^iA$. Thus the twisted coboundary of
$\delta^i\theta(y)$ is integral. Its unique primitive has norm
at most $p^{2+v_p(e-i)}$. Multiply the entire relation by this
power of $p$, lift the now integral residue by $\theta:A\to\mathcal O_C$,
and subtract $\xi^i$ times the lift from the primitive. Its degree
increases by one with the error still integral. The total power
used in degrees below $e$ is $n=2e+v_p(e!)$. In degree $e$ the
weight-zero logarithmic coefficient of the integral error is
$\delta^e p^n\alpha$, up to sign. With $|\delta|=p^{-1/(p-1)}$
we obtain
\[
                      |\alpha|\le p^{2e+v_p(e!)+2+e/(p-1)}.
\]
% END K81C3

% BEGIN K81C4
\paragraph{K81C4. An integral cocycle with a prescribed power-sum coefficient.}
Write $u=1+p$, $U=u^{-e}$, $E=[\epsilon^{1/p}]$, $\pi=E-1$,
$\mu=[\epsilon]-1=\xi\pi$, and let $\sigma$ be arithmetic
Frobenius on $K_0$. Choose a finite list of pairs $(c_i,m_i)$ with
$c_i\in\mathcal O_{K_0}$ and $m_i$ positive and prime to $p$.
Suppose $d_0\in\mathcal O_{K_0}$ solves
$d_0-\sigma d_0=\sigma(\sum_i c_i)$. For $j\ge1$ define
$d_j=(1-p^j\sigma)^{-1}\sigma(\sum_i c_i m_i^j)$;
the inverse is a convergent geometric series. There exists an
integral continuous $\mathbb Z_p(-e)$-cocycle with logarithmic
coefficient
\[
     \frac{\sum_i c_i m_i^e\zeta_p^{m_i}+p^e d_e}
                       {\log(u)\,e!\,p^e} .
\]
Here and below use the same $p$-power roots for $t$ and $E$.
We prove the integral assertion, including the period convergence.
Embed $\mathcal O_{K_0}$ in $A$ by the constant unramified Witt
vectors. Put $Q=\sum_i c_i E^{m_i}$. There is $y\in A$ with
$(1-\varphi)y=Q$ given by
\[
 y=Q+d_0+
     \sum_{n\ge0}\sum_i\sigma^{n+1}(c_i)
                         \bigl((1+\mu)^{m_i p^n}-1\bigr).
\]
Indeed $\varphi E=1+\mu$, the constant terms give the equation
for $d_0$, and the rest telescopes. The series converges in the
$(p,\mu)$-adic topology: for every fixed degree $k>0$ the
coefficients $\binom{m_i p^n}{k}$ tend to zero $p$-adically,
while high degrees are divisible by a high power of $\mu$.
The ring of integral Witt vectors of the complete valuation ring
of the tilt is complete in this topology (modulo each power of
$p$ use its finitely many Witt coordinates and the valuation
topology). The map to $B_M$ is continuous for this topology since
$\mu^M=0$ there. In $B_M$ one has $E=\zeta_p\exp(t/p)$:
both sides have the same residue, and taking their $p$-th powers
and using the finite logarithm gives equality. The Taylor expansion
of $y$ is $\sum_{j=0}^e w_jt^j$, with
\[
 w_j=\frac{\sum_i c_i m_i^j\zeta_p^{m_i}+p^j d_j}{j!p^j}.
\]
In fact the tail's $j$-th coefficient for $j>0$ is the geometric
series defining $d_j/j!$. This also verifies the expansion when
$e=0$.

We next construct $b\in B_M$ for which $h b-b\in\mathbb Z_p$
for $h\in H$ and $y+U\gamma b-b\in\mathbb Z_p$. To do so take
integral sequences $P_k,Q_k$ with $Q_0=Q$, $P_k$ fixed by $H$,
and in the ring with $\mu$ inverted the identity
\[
  (U\gamma-1)\frac{P_k}{\mu^{k+1}}
       =\frac{Q_k}{\mu^k}-p\frac{Q_{k+1}}{\mu^{k+1}}.
\]
Here is a coefficient construction. Work on one monomial $E^m$,
using $X=\mu$ and $v=((1+X)^u-1)/X$, which is a formal unit.
If $Q_k=E^m f(X)$ set $P_k=E^m P(X)$, $a=k+1$. It suffices to
solve
\[
 U(1+X)^m P(Xv)-v^aP(X)\equiv v^a X f(X)\pmod p .
\]
On $X^i$ the left side has zero coefficient in degrees $\le i$
modulo $p$, and coefficient $m$ in degree $i+1$, since
$v\equiv1\bmod(p,X^2)$ for odd $p$. Successive coefficients
therefore determine $P$ modulo $p$. Lift integrally and divide the
error by $p v^a$, with the sign chosen in the displayed recursion,
to define the next $f$. Start with $f=1$ and evaluate the series
at $\mu$ using completeness. This uses
$\gamma(E)=E(1+\mu)$ and $\gamma(\mu)=\mu v(\mu)$.
Sum with coefficients $c_i$, all fixed by $H$ and $\gamma$.

Choose $\beta_k,\eta_k\in A$ satisfying
\[
 (\varphi-\xi^{k+1})\beta_k=P_k,\qquad
 (\varphi-\xi^k)\eta_k=Q_k,\qquad \eta_0=-y .
\]
For surjectivity, modulo $p$ the polynomial $x^p-a x-b$ with
integral coefficients has an integral root in the algebraically
closed tilt. Subtract the lift, divide the error by $p$, and repeat
by $p$-adic completeness. In $W(C^\flat)$ the element $\pi$ is
invertible and $\varphi\pi=\mu$. The last recursion then says that
\[
 (U\gamma-1)\frac{\beta_k}{\pi^{k+1}}-
       \frac{\eta_k}{\pi^k}+p\frac{\eta_{k+1}}{\pi^{k+1}}
                  \quad\hbox{is fixed by }\varphi .
\]
It belongs to $\mathbb Z_p$ since each Frobenius-fixed Witt coordinate
is in $\mathbb F_p$. Likewise $(h-1)\beta_k/\pi^{k+1}$ is in
$\mathbb Z_p$. These equalities transport to $B_M$ by clearing
powers of $\pi$ and their Galois translates in $A$: all have nonzero
residue and become units in $B_M$.

Put $b=\sum_{k\ge0}p^k\beta_k/\pi^{k+1}$ in $B_M$. This
series converges in its $p$-adic lattice topology. Indeed
$\xi=\pi^{p-1}+p(1+\pi w)$ with $w\in A$, by the binomial
formula. Modulo $(\pi^a,\xi^M)$, $1+\pi w$ is a unit, and for
$n=\lceil a/(p-1)\rceil$ the binomial expansion of
$(\xi-\pi^{p-1})^{n+M-1}$ vanishes. Thus
$p^{n+M-1}\pi^{-a}\in A_M$ inside $B_M$. Since $p>2$ this gives
$p^k\pi^{-k-1}A_M\to0$, and also $p^k\eta_k/\pi^k\to0$.
The ring $A_M$ is $p$-adically complete and separated: the
successive $\xi$-graded modules are $\mathcal O_C$, and all
successive quotients are $p$-torsion free, so that the same finite
filtration is exact modulo $p^n$ and under the surjective inverse
limits. The telescoping identity now proves the assertions for $b$;
$\mathbb Z_p$ is compact, hence closed in $B_M$.

For each $j<e$ add $w_jt^j/(u^{j-e}-1)$ to $b$. These terms
are $H$-fixed, and their twisted differences at $\gamma$ are
exactly the omitted lower terms. For the resulting $b'$ set
\[
 z(g)=D_g b'+\log\chi(g)\,(w_e/\log u)t^e\quad\text{in }B_M .
\]
It is integral scalar on $H$ and $\gamma$, and it is a continuous
cocycle on $G_1$ with values initially in $B_M(-e)$. Since the
scalar lattice $\mathbb Z_p$ is closed and invariant under the
twist, the cocycle identity and density show its values lie there
on all of $G_1$. The subspace topology on this lattice is its
usual one, again by compactness and separation. This proves the
claimed integral formula.
% END K81C4

% BEGIN K81C5
\paragraph{K81C5. The full integral image at level one.}
Set $\beta=-1$ for $e=0$ and $\beta=-e-v_p((e-1)!)$ for $e>0$.
Every element of $p^\beta\mathcal O_F$ is the logarithmic
coefficient of an integral scalar cocycle. To see this write
$x=\sum_{a=1}^{p-1}b_a\zeta_p^a$ in the integral basis over
$\mathcal O_{K_0}$ obtained by multiplying the power basis by
$\zeta_p$. The power basis is integral since $\zeta_p-1$ is
an Eisenstein uniformizer over $K_0$. Solve
$(1-p^{e+1}\sigma)d=\sigma(\sum_a b_a)$ geometrically and set
$\lambda_a=b_a+p^e d$. Then
\[
     d-p^e\sigma d=\sigma\!\left(\sum_a\lambda_a\right),
     \qquad \sum_a\lambda_a\zeta_p^a+p^ed=x,
\]
using $\sum_{a=1}^{p-1}\zeta_p^a=-1$. For $e=0$ use the pairs
$(\lambda_a,a)$ and $d_0=d$ in K81C4. The output is $x/\log u$,
which ranges over $p^{-1}\mathcal O_F$. For $e>0$ use the two
pairs $(-\lambda_a/a^e,a)$ and $(\lambda_a/a^e,ua)$ for each
$a$. Their zeroth moment vanishes, so $d_0=0$ works. Their
$e$-th moment gives $d_e=(u^e-1)d$, and $\zeta_p^{ua}=\zeta_p^a$.
Thus the output is $(u^e-1)x/(\log(u)e!p^e)$. Its scalar
multiplier has valuation $\beta$ because
$v_p(u^e-1)=1+v_p(e)$ and $v_p(\log u)=1$ for odd $p$.

Let $a=e+3+v_p(e)+\lfloor e/(p-1)\rfloor$, with $v_p(0)=0$
in this expression. The bound of K81C3 gives
$p^a\alpha\in p^\beta\mathcal O_F$. For $e>0$ compare exponents
by $v_p(e!)=v_p(e)+v_p((e-1)!)$ and
$e/(p-1)\le 1+\lfloor e/(p-1)\rfloor$; at $e=0$ use $a=3$.
Subtract from $p^aZ$ the integral cocycle just constructed with
coefficient $p^a\alpha$. This gives a rational cocycle $z'$ of
coefficient zero satisfying
\[
 z'(g)\equiv p^{a+N+h}c(g)\pmod{\mathbb Z_p}
                    \qquad (g\in G_1).
\]
% END K81C5

% BEGIN K81C6
\paragraph{K81C6. The zero coefficient and the exact restriction index.}
As a $\mathbb Q_p$-vector space the continuous rational cocycles
$Z^1(F,\mathbb Q_p(-e))$ have dimension $[F:\mathbb Q_p]+1$.
Indeed rational local Tate duality gives $H^2=0$ since $e+1>0$
and $\chi(G_1)$ is infinite. The local Euler formula gives
$\dim H^1=[F:\mathbb Q_p]+\dim H^0$, and the dimension of the
coboundaries is $1-\dim H^0$. The map to $F$ taking the logarithmic
coefficient is linear by K81C3 and surjective by K81C5 and
rational scaling. Its kernel on cocycles therefore has dimension
one. If $e>0$ the nonzero cocycle $g\mapsto\chi(g)^{-e}-1$ has
coefficient zero, with primitive $1$. Thus each element of the
kernel is a rational coboundary. If $e=0$ the nonzero unramified
degree homomorphism has coefficient zero: in the completed unramified
Witt ring the difference of Frobenius and identity is surjective
by successive Artin--Schreier lifting on the residue algebraic
closure. Choose a primitive for its value at residue Frobenius
(using the residue-degree power of Frobenius). It belongs to
$B_{\rm cris}^+\subset B_{\rm dR}^+$; density gives the primitive
equation on the whole unramified quotient, as in S80J4. The
kernel consequently consists of unramified homomorphisms. In both
cases $z'$ from K81C5 becomes a coboundary over $B_{\rm cris}$.

Run K81C2--K81C5 in every coordinate of $T_e$. Combining the
rational cocycles $z'$ and their crystalline primitives in the
chosen basis gives a continuous $V$-cocycle that is finite over
$F$ and reduces to $p^{a+N+h}\phi(f)$ there. Corestrict its
class along $F/K$. Corestriction is a finite sum on cosets,
commutes with reduction and with the map to crystalline period
cohomology, and sends a coboundary to a coboundary. Hence the
corestricted rational class is finite over $K$. On the torsion
class restricted from $K$ its reduction is $[F:K]$ times the
original class, by restriction followed by corestriction. Write
$[F:K]=p^{g_K}d$ with $d$ a $p$-adic unit and divide the rational
finite lift by $d$. Equality of the torsion cohomology classes
allows an additional torsion coboundary, precisely the coboundaries
included in the cocycle definition of $H_f^1(K,V/T_e)$. We obtain
\[
        p^{N+a+h+g_K}\phi(f)\in H_f^1(K,V/T_e).
\]
Here $g_K=v_p([F:\mathbb Q_p])-v_p([K:\mathbb Q_p])$, and
$N+a+h+g_K=\mathrm{cTBY}(Y,M_0,h_0,e)$, since the two
corrections at level one are
$e+2+v_p(e)+\lfloor e/(p-1)\rfloor+1$ and $1+v_p(e+1)$.
With the Wach estimate W81T6 this proves the entire predicate
$\mathrm{TBat}$ on the stipulated level-one data for odd $p$.

In particular at the two places of H80M5 all kernel lattices have
the required $\mathrm{HypW}$, the middle terms are crystalline
on $G_1$, and $g_K=0$. The comparison of H80M3 lets one multiply
by the remaining nonnegative power to reach $c_B$. H80M4 supplies
the three rank-one clauses. Thus the original ramified
$\mathrm{admTB}$ is proved there for each good coefficient extension
and every tested row. The separate lifting estimate on a rank-$r$
row and the global lattice-index application have their own inputs.
% END K81C6

\section{The global index with all ramified places declared as rows}

% BEGIN I81G1
\paragraph{I81G1. A unit root difference on the ramified hub.}
Use the member, good localization, point and complement of H80M1.
Write $p=\ell\in\{q,d\}$, $K=E_w$, and $K_0$ for the maximal
unramified subfield of $K$. At every $w\mid p$, R79Q2 gives
$K\subset K_0(\mu_p)$ with $[K:K_0]=2$. Since
$K_0(\mu_p)/K_0$ is totally ramified and cyclic of order $p-1$,
$K(\mu_p)/K$ is totally ramified and cyclic of order $(p-1)/2$.
Lift a generator to local inertia, then to the specified global
inertia group through the place identification. For its element
$g_0$, the residues of $\chi_c(g_0)$, $1\le c\le t$, are distinct.
Indeed they equal $\chi_1(g_0)\varepsilon_p(g_0)^{-(c-1)}$
and $t<(p-1)/2$ by the host bound $p>4tr+1$.

D79G1--D79G3 apply with their original hypotheses at this localization.
The polynomial of the character corner at $g_0$ is monic of degree
$t$ over the local Hecke algebra. At the chosen point it has the
$t$ simple residue roots just found. The Hecke coefficients have
the same residues at all points. The Hensel argument in D79G4
therefore supplies the $t$ separately labelled roots at all points.
D79G5 and the construction in D79G6 now give the ordered-root cover,
the integral projectors and the unit $\delta(\pi')$ used in G79E7. The general decomposition
record also asks for residue order $p-1$. Here the order is
$(p-1)/2$. The proofs of G79E7 and G79E11--G79E13 below use
only the listed integral blocks, root basis and unit difference.
The off-support and away-from-$p$ hypotheses are supplied by D79G7.
Thus G79E11--G79E12 give, with the actual point ideal and block
cocycle modules of G79E1 and G79E7,
\[
  \operatorname{length}_O(O/I_\pi)\le
  \operatorname{length}_O C_\rho+
  \sum_{0\le j<t}\operatorname{length}_O C_j .
\]
Each intersects the global coboundaries in zero. Moreover the column
construction of G79E13 puts the restriction of every generator in
the local cut set for the corresponding block. That construction
requires the integral projector and the intertwining map on its
columns. Both have just been supplied at $q,d$.
% END I81G1

% BEGIN I81G2
\paragraph{I81G2. Normalized lengths after enlarging coefficients.}
Use the specified bridge to identify the coefficient DVR of a
localization with $O$ at $w$. Choose any good finite coefficient
extension with rings $O\subset O'$ and fields $L\subset L'$.
The given bridge field already contains all roots of unity of order
$\operatorname{lcm}(p,\mathrm{condN})$ and splits $K$. These two
properties persist under every finite coefficient enlargement. The test over
$G_K$ keeps $K$ fixed. For a lattice $T\subset V$ put $A=V/T$.
Writing $O'$ in a free $O$-basis identifies the continuous cochains
of $T'=T\otimes_O O'$, $V'=V\otimes_L L'$ and $A'=V'/T'$
with finite direct sums of the original cochains. This identification
commutes with the differential, and so with cohomology. For $V$,
being Bloch--Kato finite is the existence of a cobounding primitive
after tensoring with the specified rational crystalline period ring.
In the basis above this is exactly the corresponding assertion in
each coordinate. For $T$ take the inverse image of the rational finite
part, and for $A$ take its image. We obtain the same base change
for finite parts and for $H^1(K,A)/H_f^1(K,A)$.

Consequently the $O'$-span in this last quotient after enlargement
of any given set of original classes is its original $O$-span
tensored with $O'$. This follows by flatness, including for any
finite subset. If $d=e(O'/O)$, a finite-length module of length
$h$ has extended length $dh$, by extending a composition series
and using $\varpi O'=(\varpi')^d O'$. The same factor multiplies
$\operatorname{length}_O(O/p)$. A uniform finite bound on the
normalized lengths over $O'$ therefore applies over $O$. To see
this without any finite-generation assumption on the full span,
apply it first to finite subsets. A strictly increasing chain of
finitely generated submodules cannot have more steps than the
uniform length bound. Starting with zero and adjoining a generator
outside the current span must therefore exhaust the span.
% END I81G2

% BEGIN I81G3
\paragraph{I81G3. One integral row map for any finite list of rank cuts.}
Work over $O'$ and write $V$ for the enlarged rational rank block
$\rho/\chi_1$, with its original extended lattice $T$. By C80L1
each rank cut $z_j$ in a given finite list has the presentation
$z_j(g)=\phi_j((g-1)E_j)$ on a stable finite submodule $Y_j$
in a sum $P_j$ of actual $V_x=R_x/\chi_1$. Take
$P=\bigoplus P_j$, $Y=\bigoplus Y_j$ and
$\phi=\sum\phi_j:Y\to V/T$. Choose a nonzero nonunit $c\in O'$
killing $\phi(Y)$, possible by finite generation. Choose a full
stable lattice $\Lambda$ of $P$, using the original integral
matrices in each summand. The submodules of $\Lambda$ defined by
$c^h v\in Y$ form an ascending chain and hence stabilize, say at
$h=n$. Then
$Y\cap c^{n+1}\Lambda\subset cY$: write a vector as $c^{n+1}v$
and apply stabilization to $v$. On $L_Y=Y+c^{n+1}\Lambda$
extend $\phi$ by zero on $c^{n+1}\Lambda$. This gives a
well-defined equivariant map from a full lattice, also when
$Y$ has smaller rational rank.

In $P\oplus V$ use the stable lattice
\[
 U=\{(y,v): y\in L_Y,\quad v\bmod T=\phi(y)\}.
\]
Projection onto $L_Y$ has kernel $0\oplus T$ and splits over $O'$
by lifting a basis. Thus $U$ is finite free and spans $P\oplus V$,
and the inclusion of $T$ induces an injection
$j:V/T\longrightarrow (P\oplus V)/U$.
For $d_{j,g}=(g-1)E_j$ in the $j$th summand of $Y$,
$(d_{j,g},\phi(d_{j,g}))$ interpreted with any lift in $V$
lies in $U$. It follows that
\[
 j(z_j(g))=(-d_{j,g},0)\bmod U
           =(g-1)(-E_j,0)\bmod U .
\]
Every class in the list is therefore killed by the same map on
local cohomology. Restricting integral functionals along $T\to U$
gives a surjective map $F:U^\vee(1)\to T^\vee(1)$
by the free splitting above. The rational source embeds
equivariantly in a sum of the permitted $V_x^\vee(1)$:
$P$ is already such a sum before dualizing, and $V$ is a direct
summand of $V_\pi$ by the specified good-point splitting.
The target with its lattice is exactly the matrix-lattice dual
in the row declaration. The representations and maps are continuous
in their finite-dimensional topologies. Hence every one of these
maps lies in the quantified class defining $\mathrm{RowAt}$.
% END I81G3

% BEGIN I81G4
\paragraph{I81G4. The rank defect is bounded by this single row length.}
Suppose the maps of I81G3 have normalized row length at most
$B_w$ for all finite lists at $w$. The full row bound is sufficient.
The quotient measured by $F.\mathrm{col}$ is
\[
 D_F=H_f^1(K,T^\vee(1))/F_* H_f^1(K,U^\vee(1)).
\]
Indeed the cocycle definition quotients the finite integral cocycles
by images of finite cocycles and target integral coboundaries.
The Tate pairing of $H^1(K,V/T)$ with $H^1(K,T^\vee(1))$
annihilates $F_* H_f^1(K,U^\vee(1))$ on all the classes in the
list. Pull the pairing back along $T\to U$ by functoriality:
the first argument becomes a coboundary by I81G3. The annihilator
of the full $H_f^1(K,T^\vee(1))$ is precisely $H_f^1(K,V/T)$
by Bloch--Kato local orthogonality for the stated de Rham pair.
This is the torsion--integral finite-part comparison on the tied
period rings in the local Bloch--Kato and Tate hypotheses. The
integral dual can equivalently be computed after restriction to
$\mathbb Z_p$. Choose a generator $b$ of the inverse different
of $O'$ over $\mathbb Z_p$ and use $x\mapsto\operatorname{Tr}(bx)$.
This identifies $O'$ with its full integral trace dual and induces
the perfect pairing with the discrete quotient. Since the tested
finite subgroup is an $O'$-module, testing after this trace detects
every scalar pairing (test all its $O'$-multiples).

The span of the classes modulo the finite part thus injects in
the Pontryagin dual of $D_F$. For a finite $O'$-module duality
preserves length, by a composition series and duality on the
finite residue field. Hence this span has length at most
$B_w\operatorname{length}_{O'}(O'/p)$, regardless of the number
of generators or their individual presentations. I81G2 gives
the same normalized bound on the entire original rank cut span.
% END I81G4

% BEGIN I81G5
\paragraph{I81G5. The length lost on each character block.}
Write $A_e=(L'/O')(-e)$ for the quotient of the integral cyclotomic
line of exponent $-e$, $0\le e<t$. For one character cut take
$P,E,Y,\phi$ as in C80L1. In $P\oplus L'$ with trivial last
summand take the lattice $U=Y\oplus O'(E,1)$ and its rational
span. The sum is direct by the last coordinate, and it is stable
since $(g-1)E\in Y$. Its rational representation embeds into a sum
of the $V_x$, using the trivial direct summand of $V_\pi$. This
representation is crystalline. Its labelled jumps are contained
in the common jump sets of $V_\pi$ by strictness and S80J1.
Projection to the last coordinate is invariant and integral on $U$,
takes $(E,1)$ to $1$ and has integral kernel exactly $Y$.
Extend $\phi$ to $U$ by zero on $(E,1)$ and then $O'$-linearly
to its rational span: $A_e$ is divisible and injective over the
DVR, since a map on a nonzero principal ideal extends by dividing
its generator value. The equivariance tested by $\mathrm{TBAt}$
is only on the integral kernel, where it holds by construction.

Consequently the first clause of $\mathrm{TBAt}$ at $c_B$ puts
$p^{c_B} z$ in the finite part for every original cut. By linearity
this holds on their span at the enlarged coefficients.
H80M4 bounds its length modulo the finite part by
$(d_w+1)c_B\operatorname{length}_{O'}(O'/p)$,
where $d_w=[K:\mathbb Q_p]$. By I81G2 the same normalized length
bound holds over the original coefficients. Above $p$, finite
classes in the trivial ratio are unramified and those in the
nontrivial ratios are zero as cohomology classes, as in H80M4.
% END I81G5

% BEGIN I81G6
\paragraph{I81G6. The all-row inequality on the original Hecke ideal.}
Fix the all-row declaration of R79Q4. Suppose at each $w\mid p$
that $\mathrm{TBAt}$ holds at $c_B$ and that the maps of I81G3
satisfy the row length bound $B_w$. Write $e_O=\operatorname{length}_O(O/p)$.
Restrict $C_\rho$ simultaneously to the local quotients by $H_f^1$
at $w\mid p$. By I81G1 and I81G4 its image has length at most
$e_O\sum_w B_w$. Its kernel maps injectively to the native Selmer
group: at $p$ it is finite by definition of the kernel, away from
$S_i,Q_i,p$ inertia vanishes by D79G7 and G79E8, and G79E7
excludes global coboundaries. All these maps use the same global
lattice and point.

For the $e$th character, first restrict to inertia at $Q_i$, losing
at most $\sum_{P\in Q_i}z_P$ by G79E8, with $z_P$ as in G79E1.
The kernel vanishes on all inertia away from $p$ by D79G7. On
this kernel restrict to the local quotients by the finite parts
at $w\mid p$. Its image has length at most
$e_O c_B\sum_w(d_w+1)$ by I81G5. The new kernel for $e=0$
consists of everywhere unramified homomorphisms, hence has length
at most $e_O v_p(h_E)$ by Artin reciprocity as in G79E6.
For $0<e<t$ the new kernel is locally a coboundary above $p$,
hence has length at most $e_O u_e$ by G79E6--G79E7.
Under the original finiteness guards all these bounds are finite.
Add them in I81G1 and divide by $e_O$ using G79E1, to get
\[
\begin{split}
 v_{\rm CC}(i,\pi)\ \le\;&
 s_\pi+\sum_{0<e<t}u_e+v_p(h_E)
       +\frac{t}{e_O}\sum_{P\in Q_i}z_P\\
 &+\sum_{w\mid p}B_w+t c_B\sum_{w\mid p}(d_w+1).
\end{split}
\]
This implies the cited $\mathrm{GlobalLatticeIndexTheorem}$ at
these twisted indices at the two primes, for each good point $\pi$
and its given complement, with every place a row. In the native expression
the character coefficient is $d_w+2$, the empty column set
contributes zero to its column sums, and the additional term
$4n_Vv_\delta(i)$ is nonnegative by R79Q4. The killing inequality
bounds the auxiliary-prime correction by $p^{-2}$ on its original
guard $tr<[E(\mu_p):E]$. This proves the required implication
from the local estimates on this hub directly for its evaluation
ideal, with no global index inequality left as a premise.
% END I81G6

\section{A pointwise bound for the global rank cocycles}

% BEGIN D81S1
\paragraph{D81S1. The integral equations for the common global column module.}
Fix a twisted localization and good point at $p=q$ or $d$ as in
I81G1. Write $G=G_E$, let $O$ be the coefficient DVR with uniformizer
$\lambda$ and fraction field $L$, and let $T$ be the given rank
lattice with action $\rho/\chi_1$. Put $V=L\otimes_O T$. Use the
ordered-root cover $X'\to X_i$ and the single integral matrix family
$E_1$ from I81G1. Form the finite direct sum $P$ of one copy of
$R_{x}/\chi_1$ over $L$ for each column at each point $x'$ of the
cover, where $x$ is the image of $x'$. Let $E\in P$ be the vector
of those columns of $E_1$, in that same order, and put
\[
 d_g=gE-E,\qquad Y_0=\sum_{g\in G}O d_g\subset P .
\]
The $d_g$ belong to the common integral coordinate lattice,
so $Y_0$ is finite free over the DVR. It is stable since
$g d_h=d_{gh}-d_g$. For each native generator of $C_\rho$ the map
on the matrix family restricts, after reassembling these columns,
to an equivariant $O$-linear map $h:Y_0\to V/T$ with
$z(g)=h(d_g)$. This is the global column calculation of G79E13:
left multiplication by $R_x(g)/\chi_1(g)$ on a matrix acts on
every column, and on its image the intertwining identity is
precisely the given action on $V/T$. The module $Y_0$ and the
vector $E$ are the same for all generators. Their $O$-linear
combinations have the same presentation.

Set $M=\operatorname{Hom}_O(Y_0,T)$, $M_L=L\otimes_O M$, with
action $(g h)(y)=g h(g^{-1}y)$. Reduction identifies
$M_L/M$ equivariantly with $\operatorname{Hom}_O(Y_0,V/T)$:
lift the images of an integral basis of $Y_0$ to $V$, and observe
that two lifts differ by an element of $M$. In bases of $M$ and
its dual let $D\subset M^\vee$ be the $O$-span of all rows of the
matrices $g-1$ acting on $M$. The module $D$ is Noetherian and
has a generating list drawn from finitely many of those rows,
because each element in a finite generating set uses only finitely
many rows. Stack such rows in a matrix. Write
$\lambda^{v_1},\ldots,\lambda^{v_s}$ for its nonzero Smith
entries, up to units, with exponents $v_a\ge0$.
Integral row and column operations give this diagonal form:
move an entry of least valuation to the first position, use its
divisibility into every remaining entry to clear its row and column,
and repeat on the smaller matrix. In particular $s=0$ is allowed.

A class $u+M$ is fixed if and only if $b(u)\in O$ for every $b\in D$,
while $u\in M_L$ itself is fixed if and only if every $b(u)=0$.
In the Smith coordinates the first $s$ coordinates of a fixed
class range over $\lambda^{-v_a}O/O$ and the others over $L/O$.
Reduction of $M_L^G$ gives exactly the latter coordinates.
Consequently the quotient
\[
 Q= (M_L/M)^G/\operatorname{im}(M_L^G)
   \simeq \bigoplus_{a=1}^s O/\lambda^{v_a}O
\]
has length $\sum_a v_a$, independent of all bases and row choices.
% END D81S1

% BEGIN D81S2
\paragraph{D81S2. The Smith bound before and after simultaneous localization.}
For $h\in(M_L/M)^G$ evaluation $g\mapsto h(d_g)$ is a continuous
cocycle by $d_{gh}=d_g+g d_h$. Indeed $h$ kills a power of $\lambda$
times $Y_0$, which is open in $Y_0$ in its subspace topology:
coordinates in a basis of $Y_0$ extend to continuous linear
coordinates on $P$. The map $g\mapsto gE$ is continuous by the
good-family condition. Evaluation followed by passage to the
cohomology class therefore defines an $O$-linear map to $H^1(G,V/T)$.

Each $R_x$ is semisimple over the specified field $L$ by
$\mathrm{GoodR}$. The character twist preserves stable subspaces
and their complements. A finite direct sum of semisimple modules
is semisimple, so the stable subspace $L Y_0$ has an equivariant
projection $P\to L Y_0$ restricting to the identity there.
Thus every $h_0\in M_L^G=\operatorname{Hom}_{L[G]}(L Y_0,V)$
extends equivariantly to a map $H:P\to V$. The cocycle obtained
from its reduction has the global cobounding vector $H(E)\bmod T$,
since $H(d_g)=gH(E)-H(E)$. It follows that evaluation on cohomology
factors through the finite module $Q$ of D81S1. By I81G1 the
native $C_\rho$ meets the global coboundaries in zero, and D81S1
puts its image in the image of $Q$. Hence for
$e_O=\operatorname{length}_O(O/pO)$ one has
\[
 \frac{\operatorname{length}_O C_\rho}{e_O}
       \le \frac{1}{e_O}\sum_{a=1}^s v_a .
\]
The joint restriction of these same global classes to all
$H^1(E_w,V/T)/H_f^1(E_w,V/T)$ for $w\mid p$ is an image of
$C_\rho$, so its normalized length satisfies the same bound.
This uses the original global matrices and requires no row length
estimate. The integer on the right before normalization belongs to
the chosen point, representation family and projector. Computing
its exact value by the displayed matrix requires a certificate that the chosen
rows generate all the integral equations, as well as a basis of
$Y_0$. For an upper bound it suffices to know the full rational
rank $s$ and to exhibit a nonzero $s$-rowed minor from sampled
equations: in the Smith basis all such minors are divisible by
$\lambda^{\sum_a v_a}$, and integral changes of basis preserve
the ideal of minors. Finite generation proves existence of generating
lists. An uncertified sample of group elements gives no stopping
criterion for having obtained them, and this proof gives no estimate
of the Smith sum in terms of the family parameters.
% END D81S2

\section{Integral lifting of the ramified rows}

% BEGIN A81R1
\paragraph{A81R1. Generators of integral finite cohomology.}
On a row of H80M2 let $p=\ell$, $K=E_w$, $n=[K_0:\mathbb Q_p]$,
$e=[K:K_0]=2$ and $D=ne$. Use the specified good coefficient
extension with integers $O$, uniformizer $\lambda$ and residue
field $k$. Write $T_j$ for the actual stable lattices and $V_j$
for their rational spans. The map $F:T_1\twoheadrightarrow T_2$
is equivariant and $\dim V_2=r$. Let $P_j$ be the module of
continuous $T_j$-valued cocycles whose classes become finite in
$H^1(K,V_j)$, $B_j$ the integral coboundaries, and put
$I=F(P_1)+B_2\subset P_2$. Then $P_2/B_2$ is generated over
$O$ by at most $(D+2)r$ elements.

Here are details of the integral finiteness argument. The standard
local finiteness theorem for a finite Galois module gives finite
$H^1(K,T_2/\lambda^j)$ for each $j$. Integral $H^1(K,T_2)$
is their inverse limit. Indeed representatives of compatible
classes can be chosen successively to reduce exactly to one
another by subtracting coboundaries of lifted vectors.
If a cocycle becomes a coboundary at every level, the finite
nonempty sets of primitives have nonempty inverse limit by compactness.
Thus $H^1(K,T_2)$ is compact and separated. Its quotient by
$\lambda$ injects in $H^1(K,T_2/\lambda)$ by the multiplication
sequence. Lift a basis of that finite-dimensional quotient.
Successive approximation in powers of $\lambda$ expresses every
class as an $O$-linear combination of the lifts: the errors vanish
at each fixed finite level and the coefficient series converge in $O$.
In particular integral cohomology and its submodule $P_2/B_2$ are
finitely generated. From the same exact sequence,
$H^1(K,T_2)[\lambda]$ is a quotient of $H^0(K,T_2/\lambda)$,
of $k$-dimension at most $r$. The structure theorem over $O$
therefore gives at most $r$ torsion generators of $P_2/B_2$.

The rationalization of $P_2/B_2$ is $H^1_f(K,V_2)$.
Indeed the compact image of a continuous rational cocycle lies in
$\lambda^{-j}T_2$ for some $j$, and a rational primitive similarly
acquires integral coordinates after multiplication. By the crystalline
Bloch--Kato formula of J80H1 this finite subspace has dimension
$\dim H^0(K,V_2)+\dim(D_{\rm dR}(V_2)/\mathrm{Fil}^0)$ over the
coefficient field, at most $(D+1)r$. Adding the torsion generators
proves the assertion. Choosing cocycle representatives gives the
source's generating set modulo $B_2$, with cardinality
$g\le(D+2)r\le2(D+1)r$. This argument works in every finite
local degree with the same rational dimension estimate.
% END A81R1

% BEGIN A81R2
\paragraph{A81R2. The coarse lifting test descends with no loss.}
Let $h=H$ from H80M2 and $K_1=K_0(\mu_p)$. This is a Galois
extension of $K$ of degree $s=(p-1)/e$ prime to $p$ by R79Q2.
Crystalline periods remain a basis over the extended unramified
scalars on restriction. De Rham comparison over $K_1$ extends the
filtered space by scalars, so its labelled jumps are obtained from
the restricted labels over $K$. The coefficient extension contains
the conjugates of $K_1$ by the good-layer hypothesis. Let $U$ be
the union of source jumps over its embeddings, a subset of $[-h,h]$.
For the restricted map put
\[
\begin{gathered}
 J=\big(\{a+j:a\in U,\ 1\le j\le h\}\setminus\{0\}\big)
                  \cup\{1,\ldots,h\},\\
 w(t)=t+v_p(t!),\qquad
 E=2w(4h)+\sum_{j\in J}(1+v_p(j))+w(h).
\end{gathered}
\]
On $K_1$ the cyclotomic level-one subgroup is the whole local
Galois group, and $[K_1:K_0]=p-1$. These are exactly the native
case B conditions. Its two constant height functions are $h,h$,
so its maximum lower height is $h$ and its maximum total height
is $2h$. Inserting these in the case B lifting exponent gives $E$.

Suppose the restricted integral finite-cohomology map contains
$p^E$ times every target finite class. For $c\in P_2$, restrict
its class, choose the given integral finite lift, and corestrict
this lift to $K$. Transfer commutes with $F$ and sends integral
classes to integral classes. It also preserves rational finiteness:
the map into cohomology with $B_{\rm cris}$ coefficients commutes
with transfer by its finite sum over cosets, so a zero image
remains zero. The composition of restriction and corestriction
on the original class is multiplication by $s$. As $s\in O^\times$
we recover a finite integral lift of $p^E[c]$ over $K$.
At the cocycle level the difference is in $B_2$, as required.
Thus the sole coarse lifting test here is on the cyclotomic
field itself; it is still an integral test.

Every $j\in J$ lies in $[-h,2h]\setminus\{0\}$. Summing
$1+v_p(j)$ over the entire allowed range costs $w(h)+w(2h)$.
Legendre's formula gives $v_p(t!)\le t/(p-1)\le t/2$ since
$p$ is odd. Hence $E\le2w(4h)+2w(h)+w(2h)\le18h$.
% END A81R2

% BEGIN A81R3
\paragraph{A81R3. The invariant parameter and its two primitive branches.}
For $K$ as above let $\Delta\subset\mu_{p-1}(\mathbb Z_p)$ be
the subgroup of order $s$ fixing $K$ in $K_1$, lifted by Teichm\"uller.
On each letter of L78W13 write $X$ for the chosen cyclotomic
variable, $[d](X)=(1+X)^d-1$, and $Q=\Phi_p(1+X)$. Define
\[
 Z=\frac1s\sum_{d\in\Delta}d^{-1}[d](X),\qquad v=Z^s.
\]
The substitutions satisfy $[d]\circ[c]=[dc]$, so $Z$ transforms
under $c\in\Delta$ by $c$. Its linear coefficient is one and
its constant coefficient zero. Successive inversion makes it an
integral formal coordinate. In $O[[Z]]$ a $\Delta$-fixed series
has coefficients only at multiples of $s$, since the differences
of distinct tame roots are units. Consequently $O[[X]]^\Delta=O[[v]]$
with unique coefficients, and $v=X^s$ times a unit.

Choose a primitive root $\zeta$ and put $a=v(\zeta-1)$.
Its valuation is $s/(p-1)=1/e$ by the unit just found. The element
$a$ belongs to the $\Delta$-fixed field $K$. The orbit under
$\mathrm{Gal}(K_1/K_0)$ has size $e$. Indeed for
$b\in\mathbb F_p^\times$ the quotient
$v(\zeta^b-1)/v(\zeta-1)$ reduces to $b^s$ in the residue
field, so distinct cosets modulo $\Delta$ give distinct values.
The monic polynomial $E_0(W)$ with roots these $e$ values lies in
$O_{K_0}[W]$. Its coefficients even lie in $\mathbb Z_p$, since
the series defining $v$ is defined over $\mathbb Z_p$ and the
same set of roots is stable under absolute Galois conjugation.
Its constant term has valuation one. For each coefficient letter
take $E_i$ to be the image of $E_0$. In $O[[X]]$ every
$X-(\zeta^b-1)$ divides $E_i(v(X))$ by integral division
as in W79H4. The product divides by successive evaluation and
cancellation at the remaining distinct roots. This product is $Q$.
In the residue field $E_i(v)$ has order $es=p-1$ in $X$,
because every non-leading coefficient is integral of positive
valuation, while $Q$ reduces to $X^{p-1}$. Therefore
$E_i(v)=Q$ times a unit. The Frobenius substitution
$X\mapsto XQ$ now gives $\varphi(v)=v E_i(v)^s$ times a unit.

Finally let $u=\chi(\gamma)$ for the source's level-one generator.
It satisfies $v_p(u-1)=1$. The quotient
$([u]^*v-v)/v$ is integral because $[u](X)$ is $X$ times a unit.
It vanishes at every $\zeta^b-1$ since $u\equiv1\pmod p$ and
$v(\zeta^b-1)\ne0$. It is divisible by $Q$ by the same
evaluation argument. The constant coefficient after division by
$Q$ is $(u^s-1)/p$, a unit since $p\nmid s$. Thus
$\gamma(v)-v=v E_i(v)$ times a unit. The polynomial $E_0$
has its specified root $a\in K$ with valuation $1/e$ under
every embedding. Assembling these identical letter computations
proves all clauses of $\mathrm{wachRingRelations}$ on the very
ramified rings for $e=2$ (and the proof applies to every tame
intermediate $K_0\subset K\subset K_1$). This supplies the ring
input to the scaled letter decomposition.
% END A81R3

% BEGIN A81R4
\paragraph{A81R4. Reduction of the refined exponent to the actual letters.}
Write $S,\eta$ for the integral parameters $S_0,A$ in H80M2 and
take $h'=h=H$. In particular $1\le h'$ and the positive endpoint
of every deep interval plus $S$ is at most $h'$. The strict
filtered quotient in H80M1 puts the jumps of $V_2$ in the same
deep-or-central set as those of $V_1$. Its central graded dimension
is at most $m$ from H80M2. These give the normalized hypotheses
$\mathrm{HypN}$ with the original central width and endpoints,
including the allowed degree-two alternative. At $q$ the two
embeddings extending one $K_0$-letter have the same deep sign by
R79Q1 and R79Q3. At $d$ they have opposite signs, since the
nontrivial element of $\mathrm{Gal}(K/K_0)$ interchanges the
CM type. The dual-twist reversal in H80M2 changes both signs
simultaneously and preserves this distinction.

Define $x=h/s$, $M=\lceil4x\rceil$ and
\[
\begin{gathered}
 \rho=w(2\lceil2x\rceil-1),\quad
 \sigma=w(M-1),\quad
 \kappa=e\big(w(S)+\lfloor\log_p h\rfloor\big),\\
 y_0=\lceil\eta+\tfrac12+\log_p h\rceil,
 \qquad y_2=2\sigma+\kappa+\rho+\max(2\sigma,y_0).
\end{gathered}
\]
Use $v,E_i$ from A81R3. On letter $c$ write $J_c$ for the
source's $\mathrm{JZ}$, formed inside integral target cocycles.
Explicitly, take the span of cocycles whose cohomology class is
the Herr class of a $\Delta$-invariant Wach pair $(u,z)$.
Here $u$ is supported on $c$, $E(v)^h u$ is in the Frobenius
span of $N_{c-1}$, and both $(\gamma_K-1)u$ and $z$ lie in
$v^MN$ and are $\Delta$-invariant. The pair must satisfy the
native Wach cocycle condition. This description uses the same
integral $N=N(T_2)$ and class map as in L78W13.

Suppose the following two estimates on these objects.
First, $p^{\rho+\sigma}P_2\subset B_2+\sum_cJ_c$.
Second, for each $c$ there are at most $e(e+1)m$ cocycles
$A_c\subset P_2$ such that
$p^{\sigma+y_2}J_c\subset I+\langle A_c\rangle_O$.
The latter is needed in pure position at $q$ and mixed position
at $d$, on the actual row map. Put $A_0=\bigcup_cA_c$.
There are $n$ letters, so $|A_0|\le ne(e+1)m$.
Multiplying the first inclusion by $p^{\sigma+y_2}$, applying
the second separately to each summand, and using $B_2\subset I$
gives
\[
 p^Y P_2\subset I+\langle A_0\rangle_O,\qquad
 Y=2\rho+4\sigma+\kappa+\max(2\sigma,y_0).
\]
This implication removes the quotient, normalization, divisibility
and ring-existence premises of the native reduction on these rows.
The two displayed hypotheses are the scaled finite-cocycle
decomposition and the exceptional-generator estimate at each letter.
% END A81R4

% BEGIN A81R5
\paragraph{A81R5. The two lengths and the constant 160.}
On the data of A81R1--A81R4 suppose the case B lifting test of
A81R2 and the two scaled letter estimates of A81R4 hold.
Let $Q_2=P_2/I$, $C$ the image of $\langle A_0\rangle_O$ in
$Q_2$ and $k_O=\operatorname{length}_O(O/pO)$.
The quotient $Q_2/C$ is generated by at most $g$ elements and
killed by $p^Y$. Therefore it is a quotient of $(O/p^Y O)^g$.
The submodule $C$ is generated by at most $g_0=ne(e+1)m$
elements, each killed by $p^E$ by A81R2. It is a quotient
of $(O/p^EO)^{g_0}$. Since
$\operatorname{length}_O(O/p^aO)=a k_O$ by the filtration by
powers of $p$, additivity gives the finite natural number
\[
 m_Q=\operatorname{length}_O Q_2\le(gY+g_0 E)k_O .
\]
This is the actual native $F.\mathrm{col}$ and $k_O$ is its
specified normalizer. The representatives in A81R1 allow the
quotient by $B_2$ throughout.

For the normalized rows $x=eh/(p-1)\le\tfrac32 eh/p$. From
$w(t)\le3t/2$ and $x>0$ we get
$\rho\le6x+3/2$ and $\sigma\le6x$. Also
$\kappa\le(3/2)eS+e\log_p h$ and
$y_0\le\eta+3/2+\log_p h$. Bounding the maximum by its
two nonnegative summands gives
\[
 Y\le48x+\tfrac32 eS+(e+1)\log_p h+\eta+\tfrac92
 \le80\big(eh/p+eS+e\log_p h+1+\eta\big).
\]
With $E\le18h$, $g\le2(ne+1)r$ and $g_0=ne(e+1)m$ this proves
\[
 m_Q\le160k_O\Big((ne+1)r(eh/p+eS+e\log_p h+1+\eta)
                         +ne(e+1)m(h+1)\Big).
\]
Thus the asserted explicit local premises imply the original
normalized row inequality of H80M2 with $C_A=160$, uniformly in
the good coefficient extension.
% END A81R5

% BEGIN A81R6
\paragraph{A81R6. The denominator of a rational section on a graph row.}
Suppose in addition the rational map of the row admits a
$G_K$-equivariant section $j:V_2\to V_1$. In particular this
applies to the dual of a rational direct-summand inclusion,
twisted by $1$, with the given graph lattice on the middle space.
Define $a$ to be the least nonnegative integer with
$\lambda^a j(T_2)\subset T_1$. It exists by clearing the finitely
many denominators in the images of an $O$-basis. Applying
$\lambda^a j$ to any $c\in P_2$ gives an integral cocycle in
$P_1$: the map on $B_{\rm cris}$-cohomology carries a rational
primitive to a rational primitive. Its image under $F$ is
$\lambda^a c$ as a cocycle. In particular $F.\mathrm{col}$ is
finite on each such individual row without an analytic lifting
assumption, with
\[
 \operatorname{length}_O(P_2/I)\le (ne+2)r\,a .
\]
For an integral section it is zero. The bound applies using any
known denominator $a$ and can replace the coarse test whenever
$a\le E k_O$. It also proves A81R5's final inequality directly
whenever the displayed length bound is within that budget.
Rational splitting by itself supplies no numerical comparison of
$a$ with $E k_O$: integrality of the section is measured on the
actual lattice. For the unrestricted rows of H80M2 this argument supplies
no uniform bound on section denominators. A81R7--A81R11 below
prove the scaled target decomposition independently of a section.
% END A81R6

% BEGIN A81R7
\paragraph{A81R7. A basis and the two difference equations on a ramified letter.}
Take any crystalline $T$ with jumps in $[-h,h]$ on the present
good layer. In the Taylor expansion at $1$ its integral basis
of L78W1 becomes a basis of $D_i[[Y]]$ by W79H3. The crystalline
constants are fixed by $G_K$ by W79H1. The Taylor action
$Y\mapsto dY$ thus gives the identity matrix modulo $Y$ for
every $d\in\Delta$. It gives the identity modulo $X$ on the
original basis, since $1+X=\exp(Y)$ there. Average its elements
over $\Delta$, dividing by $s\in O^\times$. The change of basis is
the identity modulo $X$, so its determinant is a unit in $O[[X]]$.
On this invariant basis $N_i^\Delta$ is free over $O[[v]]$ by A81R3.
In particular the action of the level-one generator $\gamma$ is
the identity modulo $v$ on this module: its differences are
invariant and vanish at $X=0$. On the current fields the native
$\gamma_K=\gamma$. Every cyclotomic unit occurring on $G_K$ has
$s$th power congruent to one modulo $p$, and $\chi(\gamma)^s-1$
has valuation exactly one. The minimality and choice rule for
$\gamma_K$ therefore selects $\gamma$. This proves the module
part of $\mathrm{p9}$ at $e=2$ and the action clause
$\mathrm{GamPi}$.

Put $c_0=\chi(\gamma)^s$. On $v^jN_i^\Delta/v^{j+1}N_i^\Delta$
the action of $\gamma$ is multiplication by $c_0^j$:
$\gamma(v)/v$ is invariant with constant term $c_0$, and the
basis changes by the identity modulo $v$. For $1\le j<M$
the scalar $c_0^j-1$ is a unit times $p^{1+v_p(j)}$ by the
binomial valuation formula. For $z\in vN_i^\Delta$, multiply
by this power first at $j=1$ and subtract $(\gamma-1)k_1$
with $k_1\in vN_i^\Delta$, choosing its leading coefficient to
match the residue in degree one. At each successive $j$, multiply
the residual vector and preceding correction by the new power,
and kill the coefficient in degree $j$ by an element of
$v^jN_i^\Delta$. The final discrepancy belongs to $v^MN_i^\Delta$,
and the total exponent is $\sum_{1\le j<M}(1+v_p(j))=\sigma$.
This includes $M=1$, when no correction is needed, and gives the
native $\mathrm{GradN}$.

For $M=\lceil4h/s\rceil$, W79H4 gives $Q^h\varphi N\subset N$
and $Q^hN_i\subset O[[X]]\varphi N_{i-1}$. The second inclusion
is exactly the required $\mathrm{NinLam}$ for invariant elements.
By A81R3 the first gives
$\varphi(v^MN)\subset v^M Q^{sM-h}N$.
Since $sM\ge4h$ and $h\ge1$, this is contained in
$(\lambda,X)v^MN$. Frobenius preserves $(\lambda,X)$, so its
$j$th iterate on $v^MN$ lands in $(\lambda,X)^jv^MN$.
The finite free power-series lattices are complete and separated
for that ideal. Hence for every $b\in(v^MN)^\Delta$ the convergent
sum $-\sum_{j\ge0}\varphi^j b$ uniquely solves
$(\varphi-1)z=b$ there. Indeed its partial sums telescope,
invariance is closed and preserved by Frobenius, and a fixed
vector lies in every ideal power. This proves $\mathrm{FrobPiN}$
on the given ramified module.
% END A81R7

% BEGIN A81R8
\paragraph{A81R8. A uniform power of the Wach variable lifts across a surjection.}
On any of these level-one fields, including $K_1$, fix an actual
onto map $f:T'\to T$ of crystalline lattices whose labelled
jumps lie in $[-h,h]$. Put $H_*=2h$. On every letter
\[
                 X^{H_*}N_i(T)\subseteq f N_i(T').
\]
The argument must give an exponent independent of the depth
of reduction of $O$. Write $A_j=(O/\lambda^j)[[X]]$ for $j\ge1$,
$\epsilon=1+X$ and $\varphi(X)=XQ$. This series ring is free
over its Frobenius image on $1,\epsilon,\ldots,\epsilon^{p-1}$.
Indeed modulo $\lambda$ one has $\varphi(X)=X^p$, so
$1,X,\ldots,X^{p-1}$ is a basis by collecting degrees modulo $p$.
The stated basis is the unit triangular binomial change of this
one. For surjectivity lift an expansion modulo $\lambda$ and
expand the error successively after division by $\lambda$.
For uniqueness reduce an alleged relation modulo $\lambda$ and
repeat on the coefficient lifts. Frobenius fixes $O$ and respects
these reductions. The same basis decomposition holds after
inverting $X$: $Q$ is a unit there modulo each $\lambda^j$ since
it reduces to $X^{p-1}$. Equivalently localization at $\varphi(X)$
agrees with localization at $X$, by the nilpotent coefficient ideal.

Let $L_i$ be the image of the source Wach lattice in the target
modulo $\lambda^j$ and $I_i\subset A_j$ its annihilator ideal
on the quotient. After inverting $X$ the map is onto. To verify
this first modulo $\lambda$, compare in the Witt basis of
L78W3 on the $i$th idempotent. These bases extend to invertible
comparison matrices, and $f$ between the coefficient lattices
is onto. Thus its matrix modulo $\lambda$ has full row rank
after extending from $k((X))$ to the nonzero $i$th component
of $k\otimes_{\mathbb F_p} C^\flat$. This is a faithfully flat
extension of the field, so the Laurent matrix is onto.
The field embedding and nonzero idempotent follow as in the
coefficient argument of R80D1: the reduction of the series
variable is the nonzero tilt parameter, and an idempotent
divisible by $\lambda$ in the separated Witt coefficient ring
vanishes. The same map after Laurent localization modulo
$\lambda^j$ is onto by Nakayama with the nilpotent ideal
$(\lambda)$. Consequently $X^{b_j}\in\bigcap_i I_i$ for some
$b_j$, by taking the maximum of the finitely many powers needed
to clear preimages of the target basis elements.

For $g\in I_i$ expand
$Q^{2h}g=\sum_{t=0}^{p-1}\epsilon^t\varphi(g_t)$.
Each $g_t$ belongs to $I_{i-1}$. In fact take
$z\in N_{i-1}(T)$ modulo $\lambda^j$. Then
$Q^h\varphi z\in N_i(T)$ by W79H4, and multiplication by $g$
puts this vector in $L_i$. Another $Q^h$ puts the result in
$A_j\varphi L_{i-1}$, by the other W79H4 inclusion on the source.
Write the latter element as $\sum_t\epsilon^t\varphi(z_t)$
with $z_t\in L_{i-1}$ using the series basis decomposition.
On the other hand it equals
$\sum_t\epsilon^t\varphi(g_t z)$. The components are unique:
after inverting $X$, Frobenius of the target basis is again a
basis by the two height inclusions and invertibility of $Q$.
Compare its coordinates using the Laurent decomposition.
Localization is injective on the original free module since $X$
is regular even over the truncated coefficient ring. Thus
$g_tz=z_t$ and the assertion follows.

Set $I=\bigcap_i I_i$. The cyclic permutation of the letters
shows that the coefficients of $Q^{H_*}g$ belong to $I$
whenever $g\in I$. Therefore, if
$X^{H_*}\sum_t\epsilon^t\varphi(a_t)\in I$, all
$X^{H_*}a_t$ are in $I$: multiply by $Q^{H_*}$ and use
$Q^{H_*}X^{H_*}=\varphi(X^{H_*})$. Apply this coefficient
extraction repeatedly to
$X^{H_*}(\epsilon-1)^{b_j}\in I$. Choose $d$ with $p^d>b_j$.
At each extraction the exponent in $\epsilon^k$ loses its last
base-$p$ digit since $\epsilon^p=\varphi(\epsilon)$.
After $d$ extractions one obtains $X^{H_*}$ times each
coefficient of the polynomial $(Z-1)^{b_j}$. Its leading
coefficient is one. This proves $X^{H_*}\in I$ for every $j$.
Finally the integral cokernel is finite over the noetherian
local ring $O[[X]]$ and is $\lambda$-adically separated by the
Krull intersection theorem. The stated inclusion follows by
passing to all $j$; the initial nonuniform $b_j$ does not appear.
% END A81R8

% BEGIN A81R9
\paragraph{A81R9. From the variable annihilator to a power of the local prime.}
On $K$ or $K_1$ put $a=\lceil2h/s\rceil$, with $s=1$ on
$K_1$. Assume $h\ge1$. Fix a letter $i$, write $L=N_i(T)^\Delta$
and $L'=fN_i(T')^\Delta$, meaning the image of the invariant
vectors. Since $v$ is a unit times $X^s$, A81R8 gives
$v^a L\subseteq fN_i(T')$. Averaging the preimage over
$\Delta$ gives $v^aL\subseteq L'$. On
$v^jL/v^{j+1}L$ and on $v^jL'/v^{j+1}L'$ the operator
$g=\gamma$ acts by $c_0^j$ from A81R7, for $c_0=\chi(\gamma)^s$.
For $L'$ this follows by applying $f$ to the congruence on the
source. Both filtrations are stable under $g$.

For $z\in L$ consider $v^az$ modulo $v^a L'$. It is in $L'$,
so the polynomial $P=\prod_{j=0}^{a-1}(Z-c_0^j)$ in $g$ kills it
there, successively raising the $L'$-filtration to degree $a$.
It is also in $v^a L$, so the complementary polynomial
$R=\prod_{j=a}^{2a-1}(Z-c_0^j)$ sends it into
$v^{2a}L\subseteq v^aL'$. For $0\le k<2a$ the $p$-valuation
of the Lagrange denominator at $c_0^k$ is
\[
 2a-1+v_p(k!)+v_p((2a-1-k)!)\le w(2a-1).
\]
Indeed each difference has valuation $1+v_p(j-k)$ because
$p\nmid s$. For $d=p^{w(2a-1)}$ the polynomials
$d\prod_{j\ne k}(Z-c_0^j)/\prod_{j\ne k}(c_0^k-c_0^j)$
are therefore integral. Their sum equals $d$ by interpolation.
For $k<a$ each is a multiple of $R$, and for $k\ge a$ each
is a multiple of $P$. Thus $dv^az\in v^a L'$.
Cancel the regular scalar $v^a$ to obtain $dz\in L'$.

This proves the exact $\mathrm{WKill}$ needed in A81R4 at
$d=p^\rho$ for every onto crystalline lattice map of the
specified height, including a projection of a crystalline
extension onto the trivial line. On $K_1$, $\Delta$ is trivial,
$a=2h$ and the same argument proves that $p^{w(4h-1)}$ kills
the full Wach cokernel. In particular the native cyclotomic
case B Wach exponent $w(4h)$ is available on every map entering
the cohomological lifting construction. Passing from Wach
vectors to finite integral cocycles still needs compatible
Herr pairs and their finiteness.
% END A81R9

% BEGIN A81R10
\paragraph{A81R10. Saturation at every height-one prime of the letter ring.}
Let $0\to T_0\xrightarrow{\iota}T'\xrightarrow{f}T\to0$ be an
exact sequence of full lattices in crystalline representations.
On each letter the inclusion matrix $B:N_i(T_0)\to N_i(T')$
has maximal column rank over the residue field at every height-one
prime of $A=O[[X]]$. For the prime $(\lambda)$ this follows from
injectivity on the residual lattices (the quotient is free) and the
same Witt comparison over $k((X))$ as in A81R8. This also gives full
generic rank. The rank-zero source case is immediate.

For the other primes suppose that all maximal minors have a common
irreducible factor. By Weierstrass preparation such a factor is
associated to a distinguished polynomial, and has a root $z$ in the
algebraic open unit disk. The common zero set there is finite:
any one nonzero minor, after removal of its largest power of
$\lambda$, is a unit times a distinguished polynomial.
The matrix intertwines the two integral invertible semilinear
$\gamma$-matrices. Consequently this finite set is stable under
$z\mapsto(1+z)^{\chi(\gamma)}-1$ and its inverse. A finite
orbit implies $(1+z)^{\chi(\gamma)^b-1}=1$ for some $b>0$.
The exponent is a nonzero element of $\mathbb Z_p$; after removal
of its unit part, which acts invertibly on the open unit group,
$1+z$ has $p$-power order.

At $z=0$, W79H1--W79H3 identify the specialized map with the
injective map of crystalline spaces on that letter. At
$z=\zeta-1$ of order $p$ (meaning $\zeta$ has that order), its
Taylor lattice in either representation is
$\sum_j Y^{-j}\mathrm{Fil}^jD_{\tau}[[Y]]$ by W79H2--W79H3.
Indeed these claims identify each monomial by its filtration step
and the lattice is graded. The inclusion of filtered spaces is
strict by H80M1's intersection argument applied to this sequence.
Extend an adapted basis of the subspace to a filtration-adapted
basis of the middle space. In the resulting Taylor lattices the
images of the subspace basis with the indicated $Y$ powers are
part of a basis. Thus the map has full column rank at $Y=0$.
The map at each specialization is the actual Wach map by the
morphism identity in L78W11. At a root of order $p^j$, $j>1$,
$Q(\zeta-1)\ne0$. The Frobenius matrices from the preceding
letter become invertible on evaluation there, using both height
inclusions to clear powers of $Q$; the same holds after a finite
coefficient extension containing $\zeta$. Commutation of $B$
with Frobenius transfers its rank from $\zeta^p-1$ on the
preceding letter. Induction excludes every possible common zero.
This proves the asserted rank at height one.

It follows that the image of $N_i(T_0)$ is its rational column
space intersected with $N_i(T')$. In fact solve for the preimage
by a nonzero maximal minor over $\operatorname{Frac} A$. At every
height-one localization at least one such minor is a unit, so all
coordinates are integral there. The regular local ring $O[[X]]$
is normal, hence the intersection of these localizations is $A$.
By generic ranks and $f\iota=0$ this rational column space is exactly
the kernel of $f$. The same kernel conclusion holds after applying
the Frobenius span: its columns are the Frobenius translates of
the two bases, and the inclusion matrix is $\varphi B$.
More explicitly $A$ is free on $1,\epsilon,\ldots,\epsilon^{p-1}$
over $\varphi(A)$ by lifting the expansions of A81R8 through all
coefficient powers. Hence scalar extension by $\varphi$ preserves
the exactness just proved, while the Frobenius basis maps injectively
to the actual span by the height inclusions. We obtain the exact
kernel in the Frobenius spans as well.

The images on the localized \'etale modules are still injective
on the subobject: the residual comparison in A81R8 gives an
invertible maximal minor over the residual letter field, hence a
unit there in the complete \'etale DVR. The scalars $v$ and $E(v)$
become units there since they reduce to powers of $X$ times units.
The kernel identifications therefore give the native $\mathrm{KerLam}$.
For $v^kN$, cancel $v^k$, take the preimage in $N(T_0)$ and
multiply back. For the preimage of the Frobenius span under
$E(v)^h$, take the preimage of $E(v)^h y$ in the kernel Frobenius
span, then divide by $E(v)^h$ in the \'etale module. The
idempotent-letter and $\Delta$-invariance constraints follow by
injectivity and equivariance in each case.
% END A81R10

% BEGIN A81R11
\paragraph{A81R11. The scaled letter decomposition for every finite target cocycle.}
For $c\in P_2$ use its integral extension
$T'=T_2\oplus O e_0$, where $g e_0-e_0=c(g)$.
Rational finiteness means that a $B_{\rm cris}$ primitive exists,
so subtracting it from $e_0$ extends the fixed crystalline basis
of $V_2$ by a fixed vector above one. The extension is crystalline.
Its labelled jumps are those of $V_2$ together with zero, by strict
filtered exactness, hence lie in $[-h,h]$. Normalize the Wach
generator $e$ of the actual integral trivial representation by
$1$ on its lattice, as in WLD23 and LWR5. It is fixed by Frobenius
and the cyclotomic group. In fact a Frobenius-fixed generator
compares with a constant element of $O$ by the Witt fixed-vector
identity of LWR5, and that constant must be a unit by the comparison
basis of L78W3.

By A81R9 lift $p^\rho e_i$ to $x_i\in N_i(T')^\Delta$ on each
letter, where $e_i$ here denotes the $i$th component of $e$.
The difference $(\gamma-1)x_i$ lies in $vN_i(T')^\Delta$ by A81R7
and in the kernel by its projection. A81R10 lifts it to
$z_i\in vN_i(T_2)^\Delta$. Apply A81R7 to choose
$k_i\in vN_i(T_2)^\Delta$ such that
$z_i'=p^\sigma z_i-(\gamma-1)k_i\in v^MN_i(T_2)^\Delta$.
Replace $x_i$ by $x_i'=p^\sigma x_i-\iota k_i$. It has projection
$p^{\rho+\sigma}e_i$ and $(\gamma-1)x_i'=\iota z_i'$.
Now $\varphi x_{i-1}'-x_i'$ has zero projection and is supported
at $i$. Its product with $E(v)^h$ belongs to the middle Frobenius
span by W79H4, so A81R10 gives its preimage $w_i$ with
$E(v)^hw_i$ in the kernel Frobenius span and $w_i$ invariant.
Commutation gives
\[
        (\gamma-1)w_i=\varphi z'_{i-1}-z_i'\in(v^MN(T_2))^\Delta .
\]
By A81R7 solve uniquely
$(\varphi-1)b_i=(\gamma-1)w_i$ in $(v^MN(T_2))^\Delta$.
Thus $(w_i,b_i)$ is an actual pair of the letter $J_i$ of A81R4.
Summing and using injectivity of $\varphi-1$ on the deep module
gives $\sum_i b_i=\sum_i z'_i$.

Write $x'=\sum_i x_i'$, $d=p^{\rho+\sigma}$ and use the original
Witt comparison embedding $j_D$. The Witt vector
$b=j_D(x')-d(1\otimes e_0)$ lies in the Witt tensor with $T_2$
by the integral splitting of the underlying lattices. It satisfies
$(\varphi-1)b=j_D(\sum_i w_i)$ after identifying the kernel.
For an element $g$ whose cyclotomic image is trivial or tame,
$g$ fixes $x'$, so $-\delta b(g)=d(1\otimes c(g))$.
For a lift $g$ of the chosen procyclic generator,
$j_D(\sum_i b_i)-\delta b(g)=d(1\otimes c(g))$, by the
sum identity and $(\gamma-1)x'=\iota\sum_i z_i'$.
These are precisely the two rules determining LWR5's continuous
Herr cocycle. By its additivity the sum of the classes of
$(w_i,b_i)$ is $d[c]$. Choosing integral cocycle representatives
of the separate classes therefore proves
\[
              p^{\rho+\sigma}P_2\subseteq B_2+\sum_iJ_i .
\]
This supplies the first scaled estimate of A81R4 for the original
ramified rows, including the mixed prime. In the proof of A81R5
the remaining per-letter estimate concerns lifts through the
particular row map with at most $6m$ exceptions per letter.
% END A81R11

\section{The coarse finite-cocycle lift on the cyclotomic field}

% BEGIN B81L1
\paragraph{B81L1. Intersection with the actual kernel lattice.}
Work at $F=K_0(\mu_p)$ on the case B layer of A81R2, with $p$ odd,
$A=\prod_i O[[X]]$, $Q=\Phi_p(1+X)$, and $u=\chi(\gamma)$,
$v_p(u-1)=1$. Let $N_j=N(T_j)$ for an exact sequence of crystalline
lattices $T_0\to T_1\xrightarrow{f} T_2$. Assume the middle jumps are in
$[-h,h]$; so are both outer sets by strictness (H80M1, J80H1).
We claim $D(T_0)\cap N_1=N_0$ inside $D(T_1)$.

Recall why one can take kernels in these coordinates. In bases of
the three \'etale modules, FET22--FET25 identify the maps after
extension to $W_0=O\otimes W(C^\flat)$ with those of the original
lattices tensored with $W_0$. This base change is faithfully flat
on the \'etale coefficient ring: modulo the uniformizer of $O$ it
is, componentwise, extension of the Laurent series field into
copies of the tilt field. The base components are complete DVRs
with this uniformizer (successive lifting of a residue-field inverse
inverts every element of valuation zero); the target is torsion-free,
and its reduction is nonzero on every component. Thus exactness
and injection descend. For an element of $D(T_0)$ mapping into $N_1$, its image lies in
$N_1\cap\ker f$. A81R10 identifies this intersection with the image
of $N_0$ on each letter. By injectivity of the \'etale map its
preimage is therefore in $N_0$. Conversely functoriality carries
$N_0$ into $N_1$. In particular both $Q^{-h}N_1$ and $XN_1$
intersect the kernel in the corresponding modules for $N_0$:
multiply by $Q^h$ for the first assertion, and for the second
write the vector as $Xn$ and use injectivity of $X$ on the target.
These statements apply also to the crystalline extensions used below.
% END B81L1

% BEGIN B81L2
\paragraph{B81L2. A bounded Herr pair gives a crystalline extension.}
Call $(y,z)$ a bounded pair for $T$ if $Q^hy\in N(T)$,
$z\in XN(T)$ and $(\gamma-1)y=(\varphi-1)z$. It gives an integral
finite class. We construct the comparison for its extension to
avoid assuming that this extension is crystalline.
Write $j$ for the Witt embedding and choose $b$ in $W_0\otimes T$
with $(\varphi-1)b=j(y)$ as in LWR5. Let $c$ be the actual
cocycle constructed there: $c(g)=d(g)-(g-1)b$ inside the Witt
tensor, where $d$ is zero on the cyclotomic kernel and
$d(\gamma)=j(z)$ on a generator lift. Give $T'=T\oplus O e$
the action $g e=e+c(g)$. In its Witt tensor the vector $v=e+b$
is fixed by the kernel, $(\varphi-1)v=j(y)$, and the generator
satisfies $(\gamma-1)v=j(z)$. The columns of a Wach basis of
$N(T)$ together with $v$ form a Witt basis: its change of basis
is triangular with the last diagonal entry one. Let $N'$ be their
$A$-span. It is free. The generator and its inverse preserve $N'$
and act identically modulo $X$ since $z\in XN(T)$. These facts
hold on the full cyclotomic group by density: $N'$ and $XN'$ are
compact images of finite powers of $A$ for its $(\lambda,X)$-adic
topology in the Hausdorff Witt-coordinate topology, with continuous
substitution as in LWR3--LWR5. The kernel fixes each basis column.

The Frobenius matrix in these columns is triangular by blocks with
last diagonal entry one, invertible after inverting $Q$, and both
it and its inverse have bounded powers of $Q$ as denominators.
This follows for its old block from W79H4, for its new upper column
from $Q^hy\in N(T)$, and then for the inverse by block inversion.
These verified matrix properties suffice for the following construction.
L78W2--L78W3 clear powers of $X$ in this Witt basis and its inverse;
apply the map from $A_0$ to $L\otimes B_{\rm cris}$ and divide by
these powers there using L78W4. The two resulting matrices remain
inverses. In particular these columns can be evaluated in the
crystalline tensor space of $T'[1/p]$ before its invariant dimension
has been determined. For the gamma matrix $G$, with $G(0)=1$,
L78W9 solves $G\gamma(U)=U$ and $U(0)=1$ by dividing the $j$th
coefficient equation by $1-u^j$, $j>0$. The commuting Frobenius
matrix and its inverse have just been shown to have finite $Q$-poles,
so the continuation and evaluation in L78W10 apply to $U$ and its
inverse. At every group element uniqueness of the normalized solution
intertwines this matrix with its Galois transform. Multiplying the
columns by the evaluated $U$ thus gives a $B_{\rm cris}$-basis of
invariants, with coefficients on the same finite extension as in
W79H1. Restricting the matrices along each coefficient idempotent,
or taking all coefficient coordinates, gives the required crystalline
comparison. In particular the trivial quotient has a fixed crystalline
lift: project the invariant basis, and use that the invariant scalars
are $L\otimes K_0$ to lift the quotient generator componentwise.
Subtract this lift from $e$ to obtain a crystalline primitive for $c$.
This proves the claimed finiteness from the bounded pair itself.
% END B81L2

% BEGIN B81L3
\paragraph{B81L3. The chosen finite class has a rationally liftable bounded pair.}
Write $b_W=w(4h)$ and use the original Wach cokernel bound of
A81R8--A81R9 on the cyclotomic field. For a target integral
finite cocycle $c$ form $\widetilde T_2=T_2\oplus Oe$ using $c$.
The rational extension is crystalline: the assumed crystalline
primitive yields an invariant lift of $e$ over $B_{\rm cris}$,
which together with a comparison basis in $V_2$ supplies a basis.
By strictness the extra graded piece has degree zero. In
$N(\widetilde T_2)$ lift $p^{b_W}$ times the fixed generator of
$N(O)$ to $v_2$. The differences
$y=(\varphi-1)v_2$, $z=(\gamma-1)v_2$ lie in the kernel \'etale
module. They give a bounded pair by the lower-height bound, trivial
action modulo $X$, and B81L1 on the kernel (also divide by $X$ in
the second coordinate). In the Witt tensor write $v_2=p^{b_W}e+b$.
Its two difference equations and invariance on the cyclotomic kernel
are the defining rules in LWR5, as checked also in A81R11. Accordingly its Herr class is
$p^{b_W}[c]$, with this sign.

The rational map $H^1_f(F,V_1)\to H^1_f(F,V_2)$ is surjective.
Indeed the three crystalline spaces are exact (J80H1), so
$H^1(F,B_{\rm cris}\otimes V_0)\to H^1(F,B_{\rm cris}\otimes V_1)$
is injective by the long exact sequence. The ordinary exact
sequence through $H^1(F,V_2)$ consequently restricts to an exact
sequence through $H_f^1(F,V_1)$, retaining the three $H^0$ terms.
Applying the crystalline dimension formula and strictness of the
filtered sequence shows that the image in $H_f^1(F,V_2)$ already
has its whole dimension: the dimensions of the three quotients
by $\mathrm{Fil}^0$ add in the short exact sequence and the $H^0$
terms cancel. Hence there is a rational crystalline extension
$\widetilde V_1\to\widetilde V_2$ inducing $f:T_1\to T_2$ on
kernels and the identity on quotient lines, after adjusting a
rational lift by a coboundary.

Choose a lattice $\widetilde T_1=T_1\oplus O p^l e_1$ in it with
$l$ large: compactness clears the values of the rational extension
cocycle, and the image of the added generator is integral in
$\widetilde T_2$. Its image is a lattice $I$ of finite index in
$\widetilde T_2$, containing $p^l\widetilde T_2$ (the kernel map
is onto). The integral maps $\widetilde T_2\xrightarrow{p^l}I$
and $I\hookrightarrow\widetilde T_2$ compose to multiplication
by $p^l$. Apply the Wach cokernel bound to
$\widetilde T_1\twoheadrightarrow I$ and the Wach functor to these
maps. For some $k$ this gives $v_1\in N(\widetilde T_1)$ lying
above $p^k v_2$. Its image in the \'etale quotient line is constant,
since the map on quotient lines is rationally the identity.
B81L1 and the same two difference calculations give a bounded
pair in $D(T_1)$ lying exactly above $(p^k y,p^k z)$. The size
of $k$ will have no effect on the final bound.
% END B81L3

% BEGIN B81L4
\paragraph{B81L4. The polynomial and the two integral regulators on the kernel.}
Use $N=N_0$, $D=D(T_0)$, and write $U_1$ for the set of all
jumps of the source $V_1$ over the embeddings of the cyclotomic
field. Use the set $J$ of A81R2, so that it contains each nonzero
$a+j$ with $a\in U_1$ and $1\le j\le h$. Put
\[
 P(B)=\prod_{j\in J}(B-u^{-j}),\qquad
 C=\prod_{i=1}^{h}(u^i-1),\qquad N_Q=Q^{-h}N.
\]
Then $(\gamma-1)P(\gamma)N_Q\subset N$. Indeed the quotient
$N_Q/N$ is free over $A/(Q^h)$ in a Wach basis, up to the
trivialization by $Q^h$. It is $p$-torsion free; Weierstrass
division for the distinguished polynomial $Q^h$ makes each
coefficient quotient finite free over $O$. After inverting $p$
this quotient is the product of its Taylor expansions of order
$h$ at the distinct primitive roots (Chinese remaindering).
In the coordinate $1+X=\zeta\exp(t')$, $Q$ is $t'$ times a
unit, and $\gamma$ acts on $t'$ by $u t'$. W79H2--W79H3 give
a filtered basis $b_a$ of invariant periods, indexed with
multiplicities by the jumps at this branch, in which the local
lattice has basis $(t')^{-a}b_a$. To verify the last description,
the degree-$k$ constant coefficients in the graded expansion
lattice consist precisely of $\mathrm{Fil}^{-k}$ by W79H2;
interpolation in W79H3 separates the degrees. A basis adapted
to this filtration consequently gives the displayed powers.
Thus in $Q^{-h}N/N$ the eigenvalues at this branch are
$u^{-a-j}$ for $1\le j\le h$, with no Jordan contribution.
These $a$ are jumps of $V_1$ by strictness. The claimed
polynomial kills every branch, including eigenvalue one, and
hence kills the integral quotient.

Two more calculations take place over $A$ without inverting $p$.
One has
\[
 C XN\subset(\gamma-1)XN+X^{h+1}N,
 \qquad (1-\varphi)X^{h+1}N=X^{h+1}N.
\]
For the first, $\gamma$ acts on $X^iN/X^{i+1}N$ by $u^i$.
Thus $\prod_{i=1}^h(\gamma-u^i)$ carries $XN$ into
$X^{h+1}N$. Divide this polynomial minus its value at one by
$B-1$ and absorb its sign to give the assertion. For the second,
$Q^h\varphi N\subset N$ gives
$\varphi(X^{h+1}N)\subset QX^{h+1}N$. The ideal $(p,X)$ is
Frobenius-stable and contains $Q$. Iteration therefore pushes
into $(p,X)^kX^{h+1}N$. The geometric series for $(1-\varphi)^{-1}$
converges in the complete power-series module. Finally both $p$
and $X$ are nonzerodivisors on $N_Q/N$. For $X$ this follows
from the same finite free coefficient quotient: after passing to
$L$, its roots $\zeta-1$ are nonzero. In particular
$N\cap XN_Q=XN$.
% END B81L4

% BEGIN B81L5
\paragraph{B81L5. Correcting the integral defect with the prescribed exponent.}
In the kernel module set $d(Y,Z)=(\gamma-1)Y-(\varphi-1)Z$.
The admissible inputs for $d$ form $N_Q\oplus XN$.
Both summands are $\gamma$-stable. Indeed $\gamma(Q)/Q$ is an
integral unit of constant coefficient one, since the substitution
permutes the primitive roots and $Q(0)=p$. It follows that
$(\gamma-1)N_Q\subset XN_Q$, and
$(\varphi-1)XN\subset XN_Q$ by the height bound. Thus
the defect group can be computed inside $XN_Q$.

Suppose $\omega\in XN_Q$ and, for some $k$, one has
$p^k\omega=d(Y,Z)$ with $Y\in N_Q$, $Z\in XN$. Write
$CZ=(\gamma-1)x+z_1$ with $x\in XN$ and
$z_1\in X^{h+1}N$ using B81L4. Commutativity gives
\[
 Cp^k\omega=(\gamma-1)(CY-(\varphi-1)x)-(\varphi-1)z_1.
\]
The first parenthesis is in $N_Q$; the last term is in $N$.
Apply $P(\gamma)$: its product with $\gamma-1$ kills $N_Q/N$,
so $Cp^kP(\gamma)\omega\in N$. Since $C\ne0$ and $N_Q/N$ is
$p$-torsion free, $P(\gamma)\omega\in N$. It is also in
$XN_Q$, hence lies in $XN$. The two regulators in B81L4 now give
$CP(\gamma)\omega\in d(N_Q,XN)$: the term in $X^{h+1}N$
is a $(1-\varphi)$-image there. The polynomial
$P(B)-P(1)$ is a multiple of $B-1$ over $\mathbb Z_p$, and
$\omega\in N_Q$. Subtracting its contribution gives
$CP(1)\omega\in d(N_Q,XN)$. The valuation of $CP(1)$ is exactly
\[
                 d_U=w(h)+\sum_{j\in J}(1+v_p(j)).
\]
Indeed $0\notin J$ and $v_p(u^j-1)=1+v_p(j)$ for every nonzero
integer $j$. A $p$-adic unit rescales $CP(1)$ to $p^{d_U}$.

For the pair $(y,z)$ of B81L3, lift $p^{b_W}Q^h y$ and
$p^{b_W} z/X$ integrally to $N_1$ by the Wach cokernel estimate.
Divide the first lift by $Q^h$ and multiply the second by $X$ to
obtain $(Y,Z)$ above $p^{b_W}(y,z)$. The element $d(Y,Z)$ lies
in the \'etale kernel, and is in $XQ^{-h}N_1$ by the calculation
above. B81L1 and injectivity of $X$ put it in $XN_Q$. This
defect satisfies the torsion hypothesis just treated: subtract
from $p^k(Y,Z)$ the $p^{b_W}$-multiple of B81L3's bounded
cocycle over $p^k(y,z)$. The differences are in the kernel
$N_Q\oplus XN$ by B81L1, and their differential is $p^kd(Y,Z)$.
Hence one can subtract from $p^{d_U}(Y,Z)$ a kernel element
of $N_Q\oplus XN$ whose differential is $p^{d_U}d(Y,Z)$. The result is a cocycle
pair on $T_1$ mapping to $p^{d_U+b_W}(y,z)$ in both coordinates.

By B81L2 the lifted pair has integral finite class. Naturality of
LWR5 follows with its explicit primitive by applying the lattice
map to all the Witt-coordinate equations. Its class maps to
$p^{d_U+2b_W}[c]$ by B81L3. This is exactly the case B exponent
$E$ of A81R2 at this member and on this restricted map. The
prime-to-$p$ corestriction in A81R2 gives the original coarse
lifting test at the two ramified primes, on every tested row.
% END B81L5

\section{The exceptional directions in an individual ramified letter}

% BEGIN P81F1
\paragraph{P81F1. A lattice described by Taylor coefficients.}
Use the actual local data of H80M2, a good coefficient field $L$
with DVR $O$, the invariant parameter $v$ and the degree-two
polynomial $E(v)$ of A81R3. Fix one residue embedding $c$, called
a letter in the source. Put $h=H\ge1$ and $R=O[[v]]$.
Write $D$ for the reduction of the invariant Wach module on the
preceding letter modulo $v$. A81R7 shows that $D$ is free over
$O$ and lifts to an invariant basis. Indeed a $\Delta$-invariant
element in $vN$ equals $v$ times an invariant element, by
cancellation. This identifies $D$ with the original $L_c$ quotient.
In its rational tensor the model acts on the coefficient functions
alone. For $S\subset D$ the native $\mathrm{LamLA}(S)$ is
\[
              \Lambda(S)=E(v)^{-h}(R\otimes_O S).
\]
Here and below only the chosen letter of the product is used.
The coefficient embedding is injective and $E(v)$ is invertible
in the enlarged Laurent coefficient ring used in $\mathrm{LA}$:
$E(v)$ is a unit times $Q$ on this letter (A81R3), and $Q$
reduces to a power of the invertible cyclotomic variable in the
residue Laurent field. These facts justify cancellation there.

At a primitive root $\zeta$ on this letter the source value
$\mathrm{valAt}(y)$ of a lift $y$ of $d\in D$ depends only on
$d$. Indeed the Taylor series of $\varphi y$ is regular, with
value obtained by applying Frobenius to the constant at $1$ on
the preceding letter, by W79H3. Clearing a power of $Q$ multiplies
the first nonzero possible coefficient by the corresponding power
of its leading Taylor coefficient. Dividing back gives exactly
the source value, including when that value is zero. A multiple
of $v$ in the predecessor has zero constant at $1$, giving
well-definedness. Denote this $O$-linear map by $b_\zeta$.
Its image lies in the crystalline period component, hence in
$\mathrm{FilC}_\tau(-h)$ for each associated label $\tau$:
all jumps there are at least $-h$ and all other labels are
unrestricted in $\mathrm{FilC}_\tau$.

For $x\in\Lambda(D)$ write the numerator $E^h x$ in an
integral basis of $D$. Take its Taylor series at each associated
$\zeta$, applying $b_\zeta$ to the basis vectors. Then
$x\in\mathcal N:=\mathrm{NL}$ exactly when every degree-$i$
coefficient lies in $\mathrm{FilC}_\tau(h-i)$ for every such
pair $(\zeta,\tau)$. This is the defining representation of
$\mathrm{NLset}$ using basis lifts. Conversely any allowed finite
presentation has the same Taylor coefficients by the linearity
just proved. The tests are linear, so taking the span in the
definition of $\mathrm{NL}$ preserves them and gives the claimed
description. In particular,
\[
 E^h(R\otimes_O D)\subseteq\mathcal N\subseteq\Lambda(D).
\]
For the first inclusion the numerator is divisible by $E^{2h}$.
Its jets below degree $2h$ vanish since $E$ has a simple zero
at every tested branch. At higher degrees the required filtration
contains the full period component. Multiplication by any element
of $R$ preserves $\mathcal N$, by multiplying Taylor series and
using the decreasing filtration.

Let $\gamma=\gamma_K$ and $u=\chi(\gamma)\in 1+p\mathbb Z_p$
with $v_p(u-1)=1$. This element fixes each $\zeta$ and scales
its logarithmic Taylor variable by $u$. The ratio $\gamma E/E$
is an integral unit. In fact each branch root $a$ of $E$ is
fixed under the substitution on $v$, so $v-a$ divides
$\gamma(v)-a$ integrally; apply the inverse substitution for
equality of the principal ideals. Thus the numerator of $\gamma x$
is obtained by applying this scaling to the old numerator and
multiplying by the unit $(E/\gamma E)^h$. Each Taylor test is
preserved, also for the inverse. Since $\gamma v/v$ is a unit,
$v\mathcal N$ is stable. Finally $v\mathcal N$ is saturated in
$\Lambda(D)$ over $O$. If $\lambda x=v n$ with $n\in\mathcal N$
and $\lambda$ a uniformizer of $O$, evaluating the numerator at
$v=0$ shows $E^h x\in v(R\otimes D)$. Thus $x/v\in\Lambda(D)$.
Its Taylor tests follow by dividing those of $n$ by $\lambda$
in the $L$-linear filtration spaces. This proves saturation on
the actual integral quotient.
% END P81F1

% BEGIN P81F2
\paragraph{P81F2. The rank of the invariant quotient on the two branches.}
The original $\mathrm{IIrest}(S)$ identifies, by P81F1 and the
definition of $\mathrm{InvL}$, with the fixed submodule for
$\gamma$ on
\[
        (\Lambda(S)+v\mathcal N)/(v\mathcal N).
\]
This ambient module is torsion-free over $O$ by saturation.
It is finitely generated, since $v E^{2h}\Lambda(S)\subset
v\mathcal N$ and its successive linear-factor quotients are
finitely generated as follows.

Write $a_1,a_2$ for the two distinct branch values of $v$.
Both lie in the maximal ideal of $O$ and $E=(v-a_1)(v-a_2)$
on this good field. Filter $\Lambda(S)$ by successive
multiplication by the factors of $v E^{2h}$, beginning with
$v$, then $2h$ copies of $v-a_1$, then $2h$ copies of $v-a_2$.
At each step the quotient is generated by a copy of $S$:
evaluation at the current root identifies $R/(v-a)$ with $O$
by division of convergent integral power series (also at $a=0$).
The last term maps to zero in the ambient module above, so the
induced filtration there has graded terms of rank at most
$\operatorname{rank}_O S$.

On the term at zero the scalar on the graded quotient is $1$,
since $E(0)\ne0$ and the substitution fixes zero. At a branch
$a$ let $j$ be the number of previous occurrences of $v-a$.
The multiplier of $\gamma$ on that graded piece is $u^{j-h}$.
Indeed in $X=\zeta\exp(Y)-1$ the series $v(X)-a$ has a simple
zero at $Y=0$. This follows from $v=Z^s$ in A81R3:
$Z'(X)$ is a unit at $X=\zeta-1$, $Z(\zeta-1)\ne0$ and $s$
is prime to $p$. Replacing $Y$ by $uY$ thus multiplies its leading
coefficient by $u$. Every other factor is nonzero at $a$ and
its ratio there is one. The denominator $E^h$ accounts for
the exponent $-h$. Each factor ideal is stable, so this computation
passes to the image filtration, even when the two branch values
coincide modulo the uniformizer.

After taking fraction-field coefficients, a filtration with a
scalar acting on each successive quotient bounds the dimension
of the fixed subspace by the sum of dimensions of the grades
whose scalar is one. This follows by intersecting the filtration
with the fixed space. The cyclotomic unit $u$ has infinite order.
For each branch $0\le j<2h$ only $j=h$ contributes. Including
zero there are precisely three possible grades. The invariant
quotient is a submodule of a finite torsion-free module over the
DVR, hence finite free, and its rank is at most
$3\operatorname{rank}_O S$. This proves the full native
$\mathrm{SecondPropositionFree}$ for each such letter,
$T$ and $S$, with the exact factor $[K:K_0]+1=3$. The same
flag, with its integral divisions, supplies the ring input with
at most three factors of scalar one.
% END P81F2

% BEGIN P81F3
\paragraph{P81F3. The integral factor for the deep jumps.}
At the fixed letter there are two extensions $\tau$ of its
unramified embedding since the good field splits $K$.
The set $\Omega$ in $\mathrm{piOm1}$ is the union of at most
two integer windows. After taking absolute values each consists
of at most $S$ consecutive positive integers at most $h$, since
the endpoint $b_\tau$ is positive and $b_\tau+S\le h$.
The signs distinguish the two deep orientations and do not
affect valuations. For any such interval and any $k\ge1$
the number of multiples of $p^k$ is at most
$\lfloor S/p^k\rfloor+1$, and is zero for $p^k>h$.
Summing in $k$ bounds the sum of the valuations on the interval
by $v_p(S!)+\lfloor\log_p h\rfloor$. Since
$v_p(1-u^i)=1+v_p(|i|)$ for nonzero $i$, the product
$\Pi=\prod_{i\in\Omega}(1-u^i)$ has valuation at most
\[
            \kappa=2\big(S+v_p(S!)+\lfloor\log_p h\rfloor\big)
\]
in the normalization $v_p(p)=1$. A repeated exponent in the two
windows occurs just once in the product, so the same upper
bound applies in that case. For $S=0$ the product is one.
Thus $p^\kappa/\Pi\in O$, giving the exact integral multiplier
required by $\mathrm{KapPi}$ for each pure or mixed letter.

In particular if the first frame proposition supplies a remainder
lattice $S$ of rank at most $2m$, P81F2 bounds its invariant
quotient by $6m$ generators. Any submodule of this finite free
quotient also has at most $6m$ generators over the DVR. This
settles the rank and integral multiplier in the refined count.
The first frame decomposition, lifting its individual coordinates,
and the comparison of the lifted coordinates with Herr classes
have separate content.
% END P81F3

% BEGIN P81F4
\paragraph{P81F4. The two integral changes of basis before Frobenius.}
On the preceding letter let $N=N_{c-1}^\Delta$ over $R=O[[v]]$
with reduction $D=N/vN$. Set $s=(p-1)/2$,
$M=\lceil4h/s\rceil\ge1$, and $\sigma=M-1+v_p((M-1)!)$.
Let $A=R/(v^M)$, $N_M=N/v^MN$ and $c_0=u^s$.
By A81R7 the action of $\gamma$ on $v^jN/v^{j+1}N$ is $c_0^j$.
Over $L$, every $d\in D_L$ has a unique fixed lift to $(N_M)_L$.
Indeed the graded scalars on the positive-degree kernel differ
from one, so $\gamma-1$ is an automorphism there. Correct an
arbitrary lift by its inverse. To control the denominators use
\[
 P(Z)=p^\sigma\prod_{j=1}^{M-1}\frac{Z-c_0^j}{1-c_0^j}\in O[Z].
\]
The product of denominators has normalized valuation
$\sum_{j=1}^{M-1}(1+v_p(j))=\sigma$.
The polynomial kills $v N_M$ by raising through its graded
steps, and $(\gamma-1)N_M\subseteq vN_M$. Hence
$P(\gamma)\widetilde d$ is an integral fixed lift of $p^\sigma d$,
independent of the integral lift $\widetilde d$ of $d$. This
also works for $M=1$. Extending the unique rational fixed lifts
$A_L$-linearly gives an isomorphism
$H:A_L\otimes_L D_L\longrightarrow (N_M)_L$ since its reduction
modulo the nilpotent $v$ is the identity. Its scaled map
$p^\sigma H$ is integral.

For the inverse apply the same calculation to the $A$-linear
dual of $N_M$ with conjugate semilinear action. That dual is
free over $A$ with reduction $D^\vee$, and its positive graded
scalars are the same $c_0^j$ since the action on the reduction
is trivial. Evaluating fixed lifts of a vector and a covector
gives an element of $A_L^\gamma=L$. This equality follows from
the same graded calculation on the positive degrees over $L$
(all $c_0^j-1$ are nonzero there). Their evaluation is therefore
the constant pairing of their reductions. The fixed covectors
are precisely the coordinate functionals of $H^{-1}$, so
$p^\sigma H^{-1}$ also preserves the integral lattices. Both
scaled maps are semilinearly equivariant and natural for all
maps induced by lattice morphisms: an equivariant map sends a
fixed lift to the unique fixed lift of the image of the constant.
% END P81F4

% BEGIN P81F5
\paragraph{P81F5. The full comparison frame for a ramified letter.}
Let $R_c=O[[v]]$ on the output letter and put
$J=(v^M E(v)^{4h})$. Frobenius gives a scalar map
$A\to R_c/J$ because $\varphi(v)=v E^s$ times a unit
by A81R3 and $sM\ge4h$. The invariant part $P_c$ of the
Frobenius span is free on the Frobenius images of the preceding
invariant basis: the images generate the Frobenius span before
invariants and are independent after inverting $Q$ by the two
Wach height inclusions. Hence their invariant coefficients are
exactly in $R_c$. Multiplication by $E^h$ identifies the two
native quotients $\mathrm{QLamLA}$ and $\mathrm{QLamP}$ with,
respectively, the constant $D$ module and $P_c$ modulo $J$.
Indeed their two denominators before multiplication are
$v^M E^{3h}$ times the respective integral modules.

Extend $p^\sigma H$ and $p^\sigma H^{-1}$ from P81F4 along the
Frobenius scalar map above. In the physical module send each
basis vector to its Frobenius translate. This constructs the
required comparison maps $\alpha,\beta$ on the two quotients,
composed with the quotient projections on their native domains.
Their two composites are $p^{2\sigma}$ by P81F4, and both
are linear over the invariant coefficient ring. Frobenius
commutes with the cyclotomic action and with every induced
lattice map. Multiplication by $E^h$ introduces the same unit
$E^h/\gamma(E^h)$ on both numerator modules in the cyclotomic
action. Consequently the maps intertwine the actions and the
maps of lattices in both directions, including the relation
$\mathrm{LAover}$ on finite tensor presentations.

It remains to identify the terms in the primitive Taylor comparison.
Write an arbitrary numerator as a sum of coefficient functions
times classes $d\in D$ and choose lifts of their images under
$p^\sigma H$ in $N$. At the root $1$ on the preceding letter
each lift has constant equal to the period value of $p^\sigma d$
and is fixed by $\gamma$ modulo $v^M$. In the regular period
Taylor series of W79H3 the variable $Y$ is multiplied by $u$;
$v^M$ vanishes to order $sM$ there. Thus all coefficients of
degrees $1$ through $sM-1$ vanish, since $u^j\ne1$ in these
degrees. At a primitive root on the next letter, Frobenius
applies the constant Frobenius and replaces $Y$ by $pY$.
Consequently the image numerator has Taylor coefficients below
$4h$ exactly those of the original functions multiplied by
$p^\sigma b_\zeta(d)$, as in P81F1. Replacing a representative
of the image changes the numerator by $J P_c$, which vanishes
to order at least $4h$ there by regularity. Multiplication by
$Q^k$ shifts agreement through every degree $i<k+4h$, because
$Q$ has a simple zero there. This is precisely the remaining
coefficient identity in $\mathrm{IsFrame}$, for every permitted
clearing exponent $k$ and every presentation. For objects outside
the crystalline bounded-jump test take the zero maps on their
respective domains. All functorial requirements test bounded
crystalline objects on both sides. This gives the full
$\mathrm{FrameA}$ on these fields, simultaneously for the lattices
tested at the letter.
% END P81F5

% BEGIN P81F6
\paragraph{P81F6. From one invariant letter to a finite cohomology class.}
Write $I_j$ for the physical vectors $w$ in $E^{-h}P_c$ with
$(\gamma-1)w\in v^j N^\Delta$. Compare them modulo $v^jN^\Delta$
by intersection with that submodule. The map from the depth-$M$
quotient to the depth-one quotient is injective. Indeed if a
representative is in $vN$, then for $1\le j<M$ its term in
$v^jN^\Delta/v^{j+1}N^\Delta$ is fixed as soon as it lies in
$v^jN^\Delta$. There the scalar is $u^{sj}\ne1$ on a free
$O$-module, so the term vanishes. Induction gives depth $M$.
The image contains $p^\sigma$ times the whole depth-one quotient.
For $w\in I_1$ apply the integral cyclotomic correction of A81R7
to $(\gamma-1)w\in vN^\Delta$. It supplies $k\in vN^\Delta$
with $(\gamma-1)(p^\sigma w-k)\in v^MN^\Delta$.
Project $k$ to letter $c$, which commutes with $\gamma$. The
adjustment stays in $E^{-h}P_c$ because $N_c^\Delta\subseteq
E^{-h}P_c$ by the height inclusions.

For $w\in I_M$ solve $(\varphi-1)z=(\gamma-1)w$ in
$v^M N^\Delta$ on the full set of letters by A81R7. This
gives a Herr pair. If two such first entries differ by
$a\in v^M N^\Delta$, write $a=(\varphi-1)b$ with $b$ in that
same module. The difference of the second entries is exactly
$(\gamma-1)b$, by commutativity and injectivity of $\varphi-1$
on this module. Thus the two classes agree. This gives an
$O$-linear class map on the depth-$M$ quotient. It commutes
with morphisms by the Witt cocycle formula of LWR5, or directly
by applying the morphism to the primitive in that formula.

Its values are finite classes. Indeed after restriction to
$K_1=K_0(\mu_p)$ the cyclotomic kernel and procyclic generator
are the same as before. The Witt formula restricts to the same
pair on these two subgroups. On the actual rings use the old
Wach columns with their original height inclusions after
restriction: the action of this subgroup preserves the ideal $(X)$ and is the
identity on the module modulo $X$. We have $Q^{2h}w\in N$ because $E\sim Q$ and
$Q^hP_c\subset N$, and $z\in v^M N\subset XN$.
The bounded-pair proof B81L2, increasing its height to $2h$,
constructs a crystalline primitive over $K_1$. For the rational
reduction of the integral class, corestriction after restriction
is multiplication by $[K_1:K]$. Transfer preserves finiteness
by its cochain formula and the degree is a unit. Thus the
original class is finite over $K$. All three clauses of the
native $\mathrm{ClassMap}$ and both of $\mathrm{InvLift}$ now
hold on these lattices.
% END P81F6

\section{Coefficient decomposition and lifting for a ramified letter}

% BEGIN V81P1
\paragraph{V81P1. The integral deep and exceptional summands.}
Use an actual row of H80M2 and A81R4, with coefficient field $L$,
integers $O$, height $h=h'$, and letter $c$. Put $A=O[[v]]$,
$D(v)=E_c(v)$ from A81R3, $N_c=(e_cN(T))^\Delta$, and
$U=N_{c-1}/vN_{c-1}$. A81R7 identifies $U$ with the native
$L_c$ and makes it finite free over $O$. For each of the two
embeddings $\tau$ extending $c$, the map $U_L\to D_{\rm dR,\tau}$
taking a lifted vector to its native $\mathrm{valAt}$ and projecting
to $\tau$ is an isomorphism. In fact the preceding letter at
$X=0$ gives the standard lattice in crystalline coordinates by
W79H3. The primitive expansion of its Frobenius substitutes $pY$
in this expansion and applies the invertible crystalline Frobenius
from letter $c-1$ to $c$. Its constant coefficient is therefore
independent of the primitive root. This computes the native value:
multiplication by $Q^k$ shifts the leading coefficient by $k$ and
multiplies it by the $k$th power of the linear coefficient of $Q$,
which the definition divides out. Constants of an invariant
basis give an $L$-basis, by A81R7 and W79H3. Finally the crystalline
component projects isomorphically to each de Rham component by W79H1.
Pull the filtration back to $U_L$ by this map, writing $F_\tau^i$.

For each positive deep interval $[b_\tau,b_\tau+S]$, let
$C_\tau=F_\tau^{b_\tau}$. Its codimension is at most $m$, since
only the central jumps can lie below $b_\tau$. For each negative
deep interval $[-b_\tau-S,-b_\tau]$, put
$B_\tau=F_\tau^{1-b_\tau}$, of dimension at most $m$ by the
same central dimension bound. Let $B$ be the sum of the latter
subspaces and $C$ the intersection of the former. Split the
saturated submodule $U\cap B$ off $U$, and in its integral
complement $V_0$ split off $U_0=V_0\cap C$. These splittings
exist since the quotients are finite torsion-free over the DVR.
Write $U=U_0\oplus U_r$, placing $U\cap B$ in $U_r$.
The remaining summand of $U_r$ has rank at most
$\operatorname{codim} C$ by its map to $U_L/C$. Consequently
$\operatorname{rank} U_r\le2m$. Projection $P:U\to U_0$
kills every $B_\tau$ after inverting $p$, since $U\cap B$ spans
$B$. Lift a basis $e_j$ of $U_0$ to the preceding Wach letter.
Every such lift satisfies the native positive filtration conditions.
The construction allows either collection of signs to be empty.
% END V81P1

% BEGIN V81P2
\paragraph{V81P2. The comparison of the two integral models for a given frame.}
On this letter let $\Lambda= D(v)^{-h}A\otimes_O U$.
This is the source's $\mathrm{LamLA}$ on the full coefficient
lattice, with coefficient functions supported on $c$. At a
primitive root use $1+X=\zeta\exp(Y)$; then $v$ is a unit in
$L[[Y]]$ and $D$ is $Y$ times a unit. In fact $v-r_\tau$, where
$r_\tau=v(\zeta-1)$, is itself a uniformizer of this Taylor ring:
$D$ has a simple zero by $D\sim Q$ and the cyclotomic polynomial,
and its other linear factor is nonzero at this root.
In the constant coordinates of V81P1, the native submodule
$\mathcal N_L=\mathrm{NL}$ consists exactly of those $x\in\Lambda$
whose Laurent coefficient $x_i$ at each root lies in $F_\tau^{-i}$.
Indeed the defining numerator $D^hx$ tests its coefficient at $j$
against $F_\tau^{h-j}$. Multiplication by a unit in $L[[Y]]$ and
its inverse preserve this weighted condition, since earlier
coefficients lie in smaller filtered subspaces. The same reasoning
shows that $\mathcal N_L$ is stable under $A$ and $\gamma$.
On the physical letter write $P_c$ for the invariant Frobenius
span and $\Lambda_P=D^{-h}P_c$. For $y\in\Lambda_P$ we have
$D^{2h}y\in N_c$ by W79H4. Such a $y$ belongs to $N_c$ exactly
when all its Taylor coefficients belong to the corresponding
$F_\tau^{-i}$. W79H2--W79H4 give the filtered Laurent lattice
and detect the integral membership after clearing this power of
$Q\sim D$, on all primitive roots.

Fix any frame $(\alpha,\beta)$ as in the original propositions; P81F5 constructs one on these rows.
Its maps on the two quotients are $A$-linear, $\gamma$-equivariant,
natural for the row map, and their composites multiply by
$p^{2\sigma}$. The error modules on these quotients are
$J_L=v^M D^{3h}(A\otimes U)$ and $J_P=v^M D^{3h}P_c$.
They lie in $v^j\mathcal N_L$ and $v^jN_c$, respectively,
for $0\le j\le M$. On the left use the triple zero and
$F_\tau^{-h}=U_L$, and on the right use $D^hP_c\subset N_c$.
A81R7 also gives $N_c\subset\Lambda_P$.

Suppose $y$ represents $\alpha(x)$. The Taylor clause of the frame
says that the expansions of $y$ in period coordinates and of $p^\sigma x$
coincide in every degree $i<3h$. To see the cutoff, first multiply
by $Q^kD^h$, with $k$ clearing the physical vector as required
in the definition. The asserted numerator agreement holds in
degrees $<k+4h$. The multiplying series has leading degree $k+h$
with nonzero coefficient, so division gives the stated agreement.
For $i\ge3h$ the filtration test is automatic for both vectors.
The physical coefficients belong to the same crystalline component
by W79H2 after clearing the denominator. Thus $y\in N_c$ if and
only if $x\in\mathcal N_L$. For a representative $x$ of $\beta(y)$
use $\alpha(x)=[p^{2\sigma}y]$: the same Taylor test allows
division by this nonzero scalar in the filtered $L$-spaces.
This use of scalar division is only a rational filtration test
inside the respective integral modules where $x,y$ already lie.

The analogous implications for $v^j\mathcal N_L$ and $v^jN_c$
in their respective directions, $1\le j\le M$, follow by lifting
the unmultiplied representative, applying $A$-linearity and adding
the appropriate error module. Applying this at $j=1$ to
$\gamma-1$ proves that $\alpha$ and $\beta$ carry invariant
classes to invariant classes: explicitly these are the elements
of the respective $\Lambda$ for which the difference lies in
$v\mathcal N_L$ or $vN_c$. This establishes the needed
$\mathrm{ModelF}$ comparisons for any supplied frame.
% END V81P2

% BEGIN V81P3
\paragraph{V81P3. The first frame proposition in both positions.}
Let $w$ be an invariant physical vector at depth one, and let
$x\in\Lambda$ represent $\beta(w)$. By V81P2,
$(\gamma-1)x\in v\mathcal N_L$. At a primitive expansion
$\gamma$ acts as $Y\mapsto uY$, $u=\chi(\gamma)$.
Thus for $i\ne0$ we have $x_i\in F_\tau^{-i}$: the scalar
$u^i-1$ is nonzero, and multiplication by the Taylor unit $v$
preserves the weighted filtration. Let $\Omega$ be the union of
the integer intervals $[-b_\tau-S,-b_\tau)$ at positive labels
and $[b_\tau,b_\tau+S)$ at negative labels. They are contained
in $[-h,h]$ and exclude zero. Put
$H(Z)=\prod_{i\in\Omega}(Z-u^i)$, so $H(1)=\pi_{\Omega,1}$.

Apply the integral projection $P$ of V81P1 coefficientwise to $x$.
Write $f_j$ for the $e_j$-coordinate of $P H(\gamma)x$.
Each $D^hf_j$ is integral in $A$, as is
$D^h(\gamma-1)f_j/v$. Indeed $D^h x$ and
$D^h(\gamma-1)x/v$ are integral. Substitution by $\gamma$ takes
$v$ and $D$ to unit multiples of themselves, preserving both
integrality assertions under the polynomial and the projection.
The first unit assertion follows from the invariant coordinate of
A81R3; for the second, apply the substitution to $Q=\varphi(X)/X$
and use $\gamma(X)/X$ a unit and $D\sim Q$.

At a positive label, for $i<-b_\tau-S$ one has
$F_\tau^{-i}=0$. The remaining degrees $i<-b_\tau$ are killed
by $H(u^i)$. At a negative label, for $i<b_\tau$, $i\ne0$,
we have $x_i\in B_\tau$, which $P$ kills; the subsequent degrees
$i<b_\tau+S$ are killed by $H(u^i)$. Therefore the Taylor
expansion of $f_j$ starts in degree $-b_\tau$ at a positive
label, and at a negative label it agrees with a constant through
all degrees $<b_\tau+S$. Multiplication by $D^h$ gives exactly
both numerator tests of $\mathrm{IsCoord}$. The integral
difference established above supplies its last test.

Each $f_j\otimes e_j$ belongs to the unrestricted invariant module.
In fact the expansion of $(\gamma-1)f_j/v$ still starts in degree
$-b_\tau$ at a positive label, where $e_j\in F_\tau^{b_\tau}$.
At a negative label the constant term cancels before division, so
the expansion starts in degree $b_\tau+S$, and
$F_\tau^{-b_\tau-S}=U_L$. These give every required filtration
test. Write $x_r=(1-P)H(\gamma)x$. Its numerator belongs to
$A\otimes U_r$, and its difference lies in $v\mathcal N_L$ by
subtracting the deep summands from $H(\gamma)x$. Thus $x_r$ is
in the invariant module with coefficients restricted to $U_r$.
Finally $H(\gamma)x-H(1)x\in v\mathcal N_L$, since
$H(Z)-H(1)$ is divisible by $Z-1$ over $O$ and $v\mathcal N_L$
is stable under $\gamma$. Take $H(1)x$ as the representative
of $\beta(\pi_{\Omega,1}w)$. The preceding equality modulo
$v\mathcal N_L$, the lifts of $e_j$ and $\operatorname{rank}U_r
\le 2m$ give all clauses of $\mathrm{FirstPropositionFrame}$.
% END V81P3

% BEGIN V81P4
\paragraph{V81P4. The third slot proposition at either pure sign.}
Take the given onto row map $F:T_1\to T_2$ and a frame as above.
Let $f$ satisfy the native coordinate tests, $z\in N_{c-1}(T_2)$
satisfy the positive filtration tests, and write $\bar z$ for its
class at $v=0$. Take any representative $y$ of
$\alpha_2(f\otimes\bar z)$. We prove the original lifting
assertion with $\rho$ of A81R4.

When both intervals are positive, $f$ has Laurent order at least
$-b_\tau$ on each branch and $\bar z\in F_\tau^{b_\tau}$.
Thus $f\otimes\bar z\in\mathcal N_L$ and V81P2 puts $y$ in the
physical invariant Wach lattice $N_c(T_2)$.
By A81R9 choose $y_1\in N_c(T_1)$ with $Fy_1=p^\rho y$.
Both $y$ and $y_1$ are in their respective depth-one invariant
modules: membership in the physical denominator lattice and
triviality of $\gamma$ modulo $v$ are A81R7.

When both intervals are negative, write the expansion of $f$ at
each branch as a constant plus terms of degree at least $b_\tau+S$.
This follows by dividing the negative coordinate test by $D^h$,
whose order there is $h$. Since $\gamma$ kills the constant in
the difference, $(\gamma-1)f/v$ has the same lower bound on its
order. Its numerator after multiplication by $D^h$ is integral
by the final coordinate test. At either the source or the target
all vectors lie in $F_\tau^{-b_\tau-S}$ by the jump bound.
Therefore $f\otimes\bar z'$ belongs to the invariant coefficient
module for every $z'$ of the preceding invariant Wach letter,
in either representation.

Now use A81R9 on the preceding letter to lift $p^\rho z$ to
$z_1$. Choose $y_1$ representing $\alpha_1(f\otimes\bar z_1)$.
V81P2 puts $y_1$ and $y$ in the respective depth-one physical
invariant modules. Naturality of the given frame and its linearity
give $F y_1-p^\rho y\in J_P\subset vN_c(T_2)$.
This is the native congruence in $\mathrm{LiftsBy}$, for every
representative $y$. Together the two cases give the pure-sign
$\mathrm{ThirdPropositionSlot}$ on the ramified rows at $q$ with
its prescribed exponent.
% END V81P4

\section{Lifting a coordinate on a mixed ramified letter}

% BEGIN X81M1
\paragraph{X81M1. The two branch values and the cost of a down-branch constant.}
On a mixed letter over $d$ denote the up deep endpoint by $b$,
the down endpoint by $b_D$, and put $B=b_D+S$ for the integral
width $S$ of H80M2. Thus $\eta<b,b_D$ and $b,B\le h$.
The dual sign reversal gives endpoints $s-1$ and $s+1$ respectively
in H80M2. The two values of $s$ use the archimedean data of their
respective embeddings. We allow both possible orders of $b,B$.

Write $a_U,a_D\in O$ for the up and down branch values of $v$.
At a fixed $K_0$-embedding the roots giving one label form one
orbit of $\mathrm{Gal}(K_1/K)$, by the extension-of-embeddings
argument in W79H2. The two values are distinct and
$v_p(a_D-a_U)=1/2$: by A81R3 they each have valuation $1/2$ and
their ratio reduces to the nontrivial square-class sign $-1$,
distinct from $1$ for odd $p$. Put $D_U=v-a_U$, $D_D=v-a_D$,
$E=D_U D_D$. At each branch $D_U$ or $D_D$ respectively is a
formal local coordinate over $L$ by the simple zero in P81F2.
An integral series in $v$ has integral coefficients in its
Taylor expansion in either of these coordinates since the two
$a$'s belong to the maximal ideal.

On this letter an $\mathrm{IsCoord}$ is $f=A/E^h$, $A\in O[[v]]$.
At the up branch $A$ vanishes to order at least $h-b$; at the
down branch $A-x_0E^h$ vanishes to order at least $h+B$ for
some $x_0\in L$. Divide $A$ integrally by $D_U^{h-b}$
by successive division at $a_U$, obtaining $A'$. At the down
branch $D_U$ is a rational Taylor unit. With
$\delta=a_D-a_U$ the precise congruence is
\[
      A'\equiv x_0 D_D^h(D_D+\delta)^b\pmod{D_D^{h+B}}.
\]
If $b<B$, its coefficient at degree $h+b$ gives $x_0\in O$.
If $B\le b\le B+2\eta$, its coefficient at degree $h+B-1$
gives $x_0\binom b{B-1}\delta^{b-B+1}\in O$.
By Legendre's factorial formula the valuation of a nonzero
binomial $\binom b k$ is at most $\lfloor\log_p b\rfloor$:
each nonzero level in the difference of the three floor sums
contributes at most one. Hence in this range
$p^{y_0}x_0\in O$, where
$y_0=\lceil\eta+1/2+\log_p h\rceil$, as in A81R4.
Both congruences hold at all primitive roots inducing the same
label, since the Taylor-coordinate changes are invertible and
the single rational quotient $A/E^h$ is the same.
% END X81M1

% BEGIN X81M2
\paragraph{X81M2. Lifting in the close-branch range.}
Use the frame of P81F1 and following claims with its model
compatibility V81P2. Denote the integral constant module on
the predecessor by $D_j=N_{c-1}(T_j)^\Delta/vN_{c-1}(T_j)^\Delta$.
The period map of P81F1 embeds its rationalization into the
crystalline space on the current letter by Frobenius of the
preceding constant term. It is the same map for each primitive
root, by W79H3. Write $F_D:D_1\to D_2$ for the induced map.
For the tested predecessor vector $x$ write $d\in D_2$ for its
class. Its up value lies in $\mathrm{Fil}^{b}$. In either case
of X81M1 choose $y\le y_0$ so $o=p^y x_0\in O$.

The model tensor $(p^yf-o)d$ belongs to the target $\mathcal N_2$.
Indeed the numerator $p^y A-oE^h$ vanishes to order at least
$h-b$ at the up branch and to order at least $h+B$ at the down
branch. On the up branch every possible degree $i$ therefore
satisfies $h-i\le b$, giving the required filtration test. On the
down branch $h-i\le-B$ and the full period space lies there
since $-B$ is a lower bound on jumps. These arguments use P81F1
and the integrality of $o$.

We record separately that $fd$ is invariant modulo $v\mathcal N_2$,
without spending $p^y$ in this assertion. In fact
$g=E^h(\gamma f-f)/v$ is integral by the last $\mathrm{IsCoord}$
clause. The substitution fixes each branch. At the up branch
$f$ has order at least $-b$ and at the down branch $f-x_0$ has
order at least $B$, with $\gamma x_0=x_0$. Since $v$ is nonzero
at both branches, the orders of $g$ are at least $h-b$ and $h+B$
there, respectively. The same two filtration tests give
$gd/E^h\in\mathcal N_2$. This argument also applies to any
source vector whose up value belongs to $\mathrm{Fil}^b$.

Let $w$ represent $\alpha(fd)$ as in the tested statement; it is
invariant modulo $vN_2$ by V81P2. Choose
$n\in N_c(T_2)^\Delta$ representing $\alpha((p^yf-o)d)$ using
V81P2. Lift $p^\rho n$ to $n'\in N_c(T_1)^\Delta$ by A81R9.
These vectors satisfy the $\mathrm{InvS}$ conditions modulo $v$
by $\mathrm{NinLam}$ and $\mathrm{GamPi}$ in A81R7.
For the constant tensor $d$ use A81R9 on the preceding letter
to choose $d'\in D_1$ above $p^\rho d$. A constant model tensor
belongs to $\mathrm{InvL}$ because its difference is zero.
Choose representatives $z,z'$ of $\alpha(d),\alpha(d')$.
Both lie in the required $\mathrm{InvS}$ by V81P2.
Naturality gives $Fz'-p^\rho z\in\mathrm{JP}_2$, and linearity
gives $p^y w-n-oz\in\mathrm{JP}_2$. Here
$\mathrm{JP}_2\subseteq v^M N(T_2)^\Delta$: multiply its
Frobenius-span factor by $E^h$ to make it integral by W79H4,
and its defining multiplier is $v^ME^{3h}$. Thus
$n'+oz'$ lifts $p^{\rho+y}w$ modulo $vN(T_2)^\Delta$, as required.
% END X81M2

% BEGIN X81M3
\paragraph{X81M3. The separated branches give a filtered source lift.}
Suppose now $b>B+2\eta$. Put
$q=D_U^{h-b}D_D^{h+B}$ and $f_0=q/E^h$.
For the tested $d\in D_2$ of X81M2,
$f_0d\in\mathcal N_2$ by the same Taylor tests: the numerator
has up order $h-b$ and down order $h+B$. Use V81P2 to choose
$n\in N_c(T_2)^\Delta$ representing $\alpha(f_0d)$.
Choose $n'\in N_c(T_1)^\Delta$ above $p^\rho n$ by A81R9,
and by V81P2 choose $\nu\in\mathcal N_1$ representing
$\beta(n')$. The maps of the frame are functorial for $F$ and
$\beta\alpha=p^{2\sigma}$ in their quotients. Therefore
$F_D\nu-p^{\rho+2\sigma}f_0d$ belongs to $\mathrm{JLA}_2$,
with $F_D$ extended to coefficients. Multiplying by $E^h$ gives
an integral series $H=E^h\nu\in D_1[[v]]$ satisfying
\[
 F_D H\equiv C q d\pmod {v^M E^{4h} D_2[[v]]},
                 \qquad C=p^{\rho+2\sigma}.
\]
This uses the exact quotient defining the frame, with its factor
$E^{3h}$ before clearing the denominator.

Because $\nu\in\mathcal N_1$, its numerator coefficient of
degree $i$ lies in $\mathrm{Fil}_D^{h-i}$ at the down branch
and in $\mathrm{Fil}_U^{h-i}$ at the up branch. The assertion
holds in the coordinates $D_D,D_U$ of X81M1: an invertible
change of local variable with zero constant mixes a coefficient
only with earlier degrees, whose filtration conditions are
stronger. On $D_1\otimes L$ the down filtration above $\eta$
vanishes, since all source jumps there are central or negative
deep. Thus $H$ vanishes to order $h-\eta$ at the down branch.
By integral division in a basis of $D_1$ write
$H=D_D^{h-\eta}H_1$ with $H_1\in D_1[[v]]$.
At the up branch $D_D$ is a rational Taylor unit; division by
its power preserves the same coefficient filtration bounds.
The congruence there now gives, to degree strictly below $4h$,
\[
    F_D H_1\equiv C D_U^{h-b}(D_U+a_U-a_D)^{B+\eta}d
                    \pmod {D_U^{4h}}.
\]
Take the coefficient of $D_U^j$ for
$j=h-b+B+\eta<h-\eta$. It is an integral vector $d'\in D_1$,
and its image is $Cd$ by the monic highest coefficient on the
right. The up filtration of $d'$ is at least $h-j>\eta$.
The source has no jumps strictly between $\eta$ and $b$, so
this is $d'\in\mathrm{Fil}_U^b$. This argument uses the common
rational crystalline space for both branch tests, with the two
separate labelled filtrations, as in X81M2. In particular the
down vanishing above is a statement in that space and follows
from injectivity of its period-component map.

Choose a predecessor Wach representative for $d'$. The final
$\mathrm{IsCoord}$ argument in X81M2 applies to $fd'$:
the up vector now has filtration at least $b$, and the down
filtration again contains the full source space at $-B$.
Thus $fd'\in\mathrm{InvL}$ on the source and its forward frame
representative lies in $\mathrm{InvS}$ modulo $v$ by V81P2.
Frame naturality and $F_Dd'=Cd$ show that this representative
maps to $Cw$ modulo $\mathrm{JP}_2\subset vN(T_2)^\Delta$.
Combining with X81M2 and scaling either lift by the unused
nonnegative power yields the original mixed
$\mathrm{ThirdPropositionSlot}$ at
$\rho+\max(2\sigma,y_0)$ for all the row maps at $d$.
% END X81M3

\section{The uniform rank bound on every ramified row}

% BEGIN U81L1
\paragraph{U81L1. At most six central generators per unit of central multiplicity.}
Take the physical and coefficient quotients modulo $vN^\Delta$ and
$v\mathcal N$ respectively, restricted to invariant classes at
depth one. V81P2 shows that $\alpha,\beta$ of P81F5 give
well-defined linear maps between these quotients. In each direction
one lifts any representative, using the corresponding inclusion
for $v$ times the integral lattice in V81P2. Their composites
multiply by $p^{2\sigma}$ since the errors $J_P,J_L$ lie inside
the respective $v$-lattices. The maps commute with the row map $F$.
Work at a single chosen letter.

For any $w$ at depth $M$ on the target pass to its depth-one
class. The first frame proposition V81P3 gives an equality in
the coefficient quotient expressing $\beta(\Pi w)$ as a sum of
coordinate tensors and a remainder. Here $\Pi=\pi_{\Omega,1}$,
the coordinate tensors satisfy the individual lifting hypotheses,
and all remainders lie in the image of the \emph{same}
$\mathrm{IIrest}(S_r)$ with $\operatorname{rank} S_r\le2m$
from V81P1. Choose an integral basis of that quotient once for
the whole target letter. Its size is at most $6m$ by P81F2.
Apply $\alpha$. For each basis class write $y_b$ for its physical
image at depth one. Multiplying the equality by
$p^\kappa/\Pi\in O$ (P81F3) gives $p^{2\sigma+\kappa}w$ as a
sum of physical coordinate classes and $O$-multiples of the
$y_b$. Scalars in this equality may be left outside each
coordinate lift.

At $q$ the two branches have the same deep sign and V81P4 lifts
each coordinate after multiplication by $p^\rho$. At $d$ the
two signs are opposite and X81M2--X81M3 give the exponent
$\rho+\max(2\sigma,y_0)$ for every relative position of the
two endpoints. This also dominates the pure-sign exponent.
Consequently for the native
$y_2=2\sigma+\kappa+\rho+\max(2\sigma,y_0)$ there is a source
depth-one class $a$ and coefficients $o_b\in O$ such that
\[
                p^{y_2} w=F a+\sum_b o_b y_b
                \quad\text{in the target depth-one quotient}.
\]
Indeed multiply the preceding lifting relation by the remaining
nonnegative power of $p$.

By P81F6, $p^\sigma a$ lifts to depth $M$ on the source and
every $p^\sigma y_b$ has a fixed choice of lift at depth $M$ on
the target. The depth-$M$ map to depth one is injective, so
the multiplied equality already holds between these lifted
depth-$M$ classes. Apply the finite Herr class map of P81F6.
Choose one integral finite cocycle for each target basis lift.
The difference between $p^{\sigma+y_2}$ times a tested target
cocycle and this combination belongs to the image of an integral
finite source cocycle modulo an integral target coboundary.
This is precisely membership in $F.\mathrm{ImF}$ plus the span
of the chosen at most $6m$ exceptions. The statement passes to
the entire $J_c$ by linearity. This proves the second scaled
estimate of A81R4 with no dependence on the number of tested
cocycles.
% END U81L1

% BEGIN U81L2
\paragraph{U81L2. The normalized row length at constant 160.}
Let $F:T_1\twoheadrightarrow T_2$ be any map in the original
$\mathrm{RowAt}$ at a twisted ramified index and a good point.
H80M1--H80M2 give the crystalline and numerical hypotheses used
in A81R1--A81R5 on this very map and each good coefficient
extension. A81R11 gives the first scaled inclusion in A81R4 for
every target finite cocycle. U81L1 gives its second inclusion on
each of the $n=[K_0:\mathbb Q_p]$ letters. B81L5 and A81R2 give
the coarse exponent $E$ for all target classes, including the
exceptions. Thus the length proof in A81R5 applies without an
extra local lifting premise. Its two quotients are killed by
$p^Y$ and $p^E$ and have at most $(2n+2)r$ and $6nm$ generators
respectively. Dividing their combined length by the length of
$O/(p)$ gives the asserted estimate with constant $160$.
In particular every map is covered by the rational ceiling
$B_w$ displayed in H80M2. This proves the original
$\mathrm{LengthBoundAtPlaces}$ for the all-row declaration on
the chosen hub. It includes every graph lattice of I81G3 and the
actual matrix-lattice target at $w$. The conclusion here uses
the signed windows determined by the two distinguished primes.
% END U81L2

\section{The residual rank estimate at the two distinguished primes}

% BEGIN J81H1
\paragraph{J81H1. A bound on all character blocks without a rank lifting hypothesis.}
Keep the native family hypotheses and the common hub of J80H2.
At a twisted index with $p=q$ or $d$ and a good point $\pi$,
let $C_\rho,C_0,\ldots,C_{t-1}$ be the global cocycle modules
of I81G1. Write $h_E$ for the class number and $u_e$ for the
native length of the $e$th positive twist.
The selected point satisfies the killing and finiteness guards of
J78H2. Put $Q_0=f^2t^2(r+t)$ and $u=\log b$.
With $e_O=\operatorname{length}_O(O/p)$, one has
\[
 \frac{1}{e_O}\sum_{0\le e<t}\operatorname{length}_O C_e
 \ \le\ \sum_{0<e<t}u_e+v_p(h_E)+p^{-2}+180Q_0u .
\]
Indeed K81C6 and H80M4 give all of $\mathrm{TBAt}$ at the joint
$c_B$ of H80M3, for every required good coefficient extension.
The character calculation in I81G5--I81G6 used no rank row estimate.
Its error is at most $t c_B\sum_w(d_w+1)$. By R79Q4,
$\sum_w d_w=2f$ and the number of places is at most $f$.
Insert $c_B\le60tf(r+t)u$. The auxiliary-prime term costs at
most $p^{-2}$ by the same strict degree guard and killing condition
as in I81G6.

In particular let $D_\rho$ denote the image of the simultaneous
restriction of $C_\rho$ in
the direct sum over $w\mid p$ of the local quotients by their
finite parts. Put $\delta_\rho=\operatorname{length}_O D_\rho/e_O$.
D81S2 shows this is a finite rational number for each point.
The kernel of this restriction injects in the Selmer group by
I81G6's kernel argument, which also holds without a row bound.
Consequently
\[
 v_{\rm CC}(i,\pi)\le s_\pi+\sum_{0<e<t}u_e+v_p(h_E)
                 +p^{-2}+180Q_0u+\delta_\rho.
\]
There is only one joint image here. D81S2 also gives a
point-dependent upper bound for it in terms of elementary divisors
of the global action on the common column module.
% END J81H1

% BEGIN J81H2
\paragraph{J81H2. The ramified-prime field on the original family hub.}
Let $S=S_c$ and $B=b/p+S+t+\eta+u$ as in H80M5. To obtain the
ramified field on the hub it suffices, at every guarded point of
every sufficiently late member, to prove
\[
                 \delta_\rho\le2080 f(rB+mb)+f .
\]
For this reduction, apply J81H1 and use $u>1$ and $Q_0\ge f$. The polynomial
$181 f^2t^2(r+t)$ times $u$ absorbs $180Q_0u+f$. The polynomial
has coefficient sum $362$ and total degree $5$. Apply the comparison
of J78H2 at $A=2080$ with the same constant already chosen for J80H2.
The argument tests only the point and representation required by the
original ramified field.

U81L2 proves the pointwise row bound of H80M2 with constant
$160$ for every row. In particular its bound holds for each of
the maps of I81G3. I81G4 and subadditivity over the places therefore
give $\delta_\rho\le\sum_w B_w$. H80M5's summation yields
$\sum_w B_w\le13\cdot160 f(rB+mb)+f$, as required. This
proves the displayed bound on the same chosen family hub. It gives
the native $\mathrm{Concl.ram}$ there together with the twelve
fields of J80H2. The common constant can be retained, increasing
the late-member cutoff to include the finitely many parameter
inequalities used above. The two remaining fields in the denominator
family conclusion are $\mathrm{mid}$ and $\mathrm{two}$. The complete
global denominator contract additionally asks for the original family
construction and the final bound on its $\eta$-sum.
% END J81H2

\section{The character condition throughout the middle range}

% BEGIN M81B1
\paragraph{M81B1. The ambient jump window at every middle-range place.}
Use the original family hub $X$ of F78D2. Fix a twisted index $i$,
$p=\ell_i$ outside the two type primes, in the native range
$\mathrm{RangeAD}(L,X.\mathrm{exc},p)$. Thus $p$ is outside the
host small set, $p\ne3,q,d$, and either $p\le\ell_\star$ or
$p\in X.\mathrm{exc}$. Take any $\pi\in Z_{\rm good}(i)$
and any $w\mid p$. These are the simultaneous guards for the
split and inert branches of $\mathrm{StarAD}$.

The conductor criterion gives $p\nmid N$ and $p\ne2$.
The argument of S79P1 applies to this $w$ regardless of the
declaration: the host embeds in the conductor-$N$ cyclotomic field,
unramified at $p$. Its completed bridge $K$ is therefore unramified,
$K=K_0$, and $d_w=[K:\mathbb Q_p]\le [E:\mathbb Q]=4n$.
The bridge inertia is $I_K$ by transport of the valuation and its
residue action. Every good coefficient extension $L'$ tested in
$\mathrm{admTB}$ is good at level one: $K\subset K_0(\mu_p)$,
and the given splitting of the higher cyclotomic layer implies
the same at level one by restriction of embeddings.

Put $H=\lfloor2b\rfloor$, $T=(t-1)/2$ and $j=j_i$.
At a label $\tau:K\to L'$ use the away-from-type reading of
$\mathrm{AdmJ}(L)$ on $\mathrm{GoodR}(X.r_i)$. Before removing
any summands, subtraction of the $\chi_1$-weight in S79P3 gives
the cloud $T+v+v_a$, the lows $T+v-b_a$ and the string
$0,\ldots,t-1$, with $v=e_{j-1}$ at the archimedean place.
At the conjugate orientation reflect all entries by $z\mapsto2T-z$.
The original dictionary applies over the specified coefficient
extension, as in S80J1--S80J2, so these lists cover the jumps of the actual twisted
ambient representation $R_\pi\chi_1^{-1}$. They contain at most
$(k-1)+m+t=n_V$ distinct integers, counted together.
The parameter bounds give $0\le v\le b+\eta+1$,
$0\le v_a\le S_c$ and $|v-b_a|\le\eta$ (G79E3).
Hence both signs of the cloud and lows have absolute value at most
$b+\eta+1+S_c+T$. The same bound contains the string.
Since $S_c+\eta+t+2\le b$ and $T\le t$, every integral
entry has $|z|+1\le H$ and $t\le H$. For the kernel lattices
tested below take the union of these integer lists over all
$d_w$ labels. The graded filtration decomposes over
$L'\otimes K=\prod_\tau L'$, so a jump of the total filtration
must occur on at least one label. Thus the union bounds the
total jump set as well, with cardinality at most $4nn_V$.
% END M81B1

% BEGIN M81B2
\paragraph{M81B2. The kernel exponent on the actual unramified bridge.}
Put $N=4nn_V$, $P=\lfloor\log_p(2H)\rfloor$. Since $K=K_0$,
the translated cyclotomic polynomials are Eisenstein over $K$.
The good level-one field has degree $p-1$ over $K$, and the
level-$j$ field has degree $p^{j-1}$ over the level-one field
for $j\ge1$. Hence $g_K=0$ and $\mathrm{vdegMs}(j)=j-1$.
For each crystalline $Y$ with its labelled jumps among those in
M81B1, take $M_0=h_0=H$. The level
$m_0=3+\min\{v\ge0:2H\le p^v(p-1)\}$ is at most $4+P$.
For $0\le i'\le e<t\le H$ and each jump $u\ne i'$, the
nonzero integer $u-i'$ has absolute value at most $2H$, so
$v_p(u-i')\le P$. Each inner sum in $\mathrm{cLL}$ costs at
most $N(4+2P)$. The initial degree costs $m_0-1$, and the
$e+1\le t$ terms add two to the maximum of one and their
respective inner sum. The remaining $\mathrm{cTB}$ corrections
are $e+2+v_p(e)+\lfloor e/(p-1)\rfloor+1+1+v_p(e+1)$,
at most $4t+3$ with the native $v_p(0)=0$. Consequently
\[
 \mathrm{cTBY}(Y,H,H,e)\le
 (3+P)+t\{3+N(4+2P)\}+4t+3=X.c_B(i).
\]
The equality is the defining $\mathrm{cBof}$ of DFD8 and C78B3
on this hub. Together with M81B1 this gives
$\mathrm{JointConditionTameBound}$ at its stored exponent.

For the first clause of $\mathrm{admTB}$ take the tested
crystalline representation $\widetilde Y$, fixed quotient
functional and integral lift $y$ of its generator. The
kernel has lattice $T=\ker(\pi_0)\cap\widetilde T$,
where $\pi_0$ denotes that functional. Its quotient lattice is
the trivial integral line because $\pi_0(y)=1$ and the functional
is integral on $\widetilde T$. Strict crystalline subobject
comparison (H80M1) puts the kernel's jumps among the tested
ambient jumps and makes the kernel crystalline. W79H1 and W79H5
give crystallinity over the good level-one subgroup and the
required $\mathrm{HypW}$ at heights $H,H$ on the chosen
comparison. The actual integral connecting cocycle is $gy-y$
viewed in $T$. W81T6 and K81C6 applied to it show that its
image under every tested equivariant map to the inverse
cyclotomic divisible line becomes finite after multiplication
by $p^{\mathrm{cTBY}(Y,H,H,e)}$. Multiplication by the remaining
power gives the first clause at $X.c_B(i)$.
% END M81B2

% BEGIN M81B3
\paragraph{M81B3. The possible cyclic obstruction and all remaining character clauses.}
Write $O=\mathcal O_{L'}$, $V=L'(-e)$, let $T_e$ be its given
stable lattice, and write $A_e=V/T_e$. Every allowed realization
is of this form by an integral choice of basis of the line.
For every $e\ge0$, Tate duality identifies $H^2(K,T_e)$ with
the Pontryagin dual of the invariants of the divisible twist
of weight $e+1$. Here the trace-dual coefficient lattice contributes
the inverse different, a free rank-one fractional $O$-ideal.
The cyclotomic image contains $1+p$. Its $(e+1)$st power differs
from one, so the invariants lie in the kernel of a nonzero scalar
on one copy of $L'/O$. This kernel and all its submodules are
finite cyclic over the DVR. The Pontryagin dual of $O/I$ is
also cyclic: using the trace-dual lattice $O^\vee$, it is
$I^{-1}O^\vee/O^\vee$ for $I\ne0$. This includes the zero
module when $I=O$. Thus the possible integral $H^2$ contributes
at most one cyclic summand. Rational $H^2(K,V)$ vanishes by
duality since the positive dual twist has zero rational invariants.

The local Euler formula gives
$\dim_{L'}H^1(K,V)=d_w+\mathbf1_{e=0}$.
For these crystalline lines the Bloch--Kato dimension is
$\mathbf1_{e=0}$: their de Rham filtration has full degree-zero
step, and only $e=0$ has a fixed vector.
For $e=0$ the unramified homomorphisms give the entire rational
finite line, since their Witt Frobenius primitive is crystalline
(K81C6). Every unramified homomorphism to $A_0$ lifts rationally
by lifting its Frobenius value and composing with the projection
of the residue Galois group to $\mathbb Z_p$.
Thus finite torsion cocycles for $e=0$ are precisely unramified
ones up to coboundaries. For $e>0$ they are precisely coboundaries,
since the rational finite space is zero. These are the two
native cocycle assertions, using the inertia identification of M81B1.

The image of $H^1(K,T_e)$ in $H^1(K,V)$ is a full $O$-lattice.
Indeed it spans since a continuous rational cocycle has bounded
image on the compact group. It is compact: the integral
cohomology is the inverse limit of the cohomologies of the finite
lattice quotients, which are finite by local finiteness, and the
degree-zero inverse system also consists of finite groups, so
there is no derived-limit term. For an $O$-uniformizer $\lambda$,
the quotient $H^1(K,T_e)/\lambda H^1(K,T_e)$ injects into the
finite group $H^1(K,T_e/\lambda T_e)$. Lift finitely many
generators of that quotient. Successive
approximation with these lifts converges in the compact inverse
limit, since powers of the uniformizer kill each fixed finite
lattice quotient. This gives finite generation over $O$.
Its projected image in
$H^1(K,V)/H^1_f(K,V)$ is likewise a full lattice.
The long exact sequence for $T_e\subset V$, followed by dividing
by the image of the rational finite part, consequently gives
\[
 0\longrightarrow (L'/O)^{d_w}\longrightarrow
 H^1(K,A_e)/H_f^1(K,A_e)\longrightarrow H^2(K,T_e)
 \longrightarrow 0 .
\]
Indeed the middle map before quotient has kernel the rational
cohomology modulo the integral lattice and cokernel zero by
rational $H^2=0$. Taking $p^a$-torsion is left exact. Its length
in the left term is $d_wa\operatorname{length}_O(O/p)$,
and in the right term it is at most
$a\operatorname{length}_O(O/p)$ by cyclicity.
The same bound $(d_w+1)a\operatorname{length}_O(O/p)$ therefore
holds in the middle. At $a=X.c_B(i)$, the quotient of cocycles
whose multiple is finite by the finite cocycles injects in
this torsion group, since all coboundaries are already finite.
This gives the native length clause even when $p-1$ divides $e+1$.

Together M81B1--M81B3 prove $\mathrm{admTB}$ at every tested
middle-range point on this hub, for both the split and the inert
case, also on restriction to the rows of a split declaration.
The rank-representation row estimates and the split-column
conditions in $\mathrm{StepAD}$ require their own arguments.
% END M81B3

\section{The unramified lifting exponent on a split middle-range declaration}

% BEGIN L81S1
\paragraph{L81S1. The place tests, coefficient extension and height.}
Use F78D2's hub $X$. Fix a twisted $i$ with $p=\ell_i$ outside the two
type primes, in $\mathrm{RangeAD}(L,X.\mathrm{exc},p)$ and split in the
host classification. At a column the matrix family is
$R_x\chi_1^{-1}$, for every $x\in X_i$; at a row it is
$(R_x\chi_1^{-1})^\vee(1)$. Both use the same stored integer
\[
 h=\lfloor h_B\rfloor,\quad h_B=b+S_c+\eta+t+2,\qquad
 e_S=X.e_S(i)=\left\lceil23\left(2+2h/p+
                         \log_p(2h+1)\right)\right\rceil .
\]
Indeed the canonical row data set $e_S=e_A(L,23,p)$ at split
primes. The general column test takes a function of the place
$w$. The split declaration uses the constant function $e_S$.

Fix any $w\mid p$ and any tested good coefficient extension $L'$
with integers $O$. M81B1 gives the unramified field $K=K_0$,
the good level-one layer and odd $p$. Away from the type primes,
the hub facts give crystallinity of every $R_x$ as in S80J3.
For the first character one may use the rank-one construction
even if there is no good point at $i$. The block dictionary makes
the infinitesimal character of the first Hecke character integral
after the specified shift (its entry is the negative first jump).
The rank-one case of the given reciprocity field produces its
continuous Galois character with the prescribed Weil--Deligne
parameters. Its Frobenius values agree with those of $\chi_1$
outside the finitely many excluded places, so Chebotarev and
continuity identify the two characters. At $w\mid p$ the character
is de Rham, and its Weil--Deligne parameter is unramified by the
hub's $\mathrm{MuUnr}$ condition and unramified character twist.
The bridge's de Rham/unramified Weil--Deligne criterion gives
crystallinity. Extending coefficients, dualizing and tensoring
crystalline bases gives both crystalline matrix families over $L'$.

S80J1--S80J2 give $|z|\le h_J=b+S_c+\eta+t$ for every jump
of the twisted family; the dual twists have jumps $-1-z$.
Both lists therefore lie in $[-h,h]$. A test in $\mathrm{LiftAt}$
uses an onto lattice map $f:T_1\to T_2$ with rational representations
subquotients of finite sums from the indicated family. Crystalline
comparison and strictness on subquotients (H80M1) give
crystallinity and the same jump window for the actual lattices
and their kernel. The following proof is uniform in their ranks.
% END L81S1

% BEGIN L81S2
\paragraph{L81S2. Tame invariants and the two integral Wach lifts.}
Put $s=p-1$ and $\Delta=\mu_{p-1}(\mathbb Z_p)$, acting by
$[d](X)=(1+X)^d-1$. Apply A81R3's integral coordinate
$Z=s^{-1}\sum_{d\in\Delta}d^{-1}[d](X)$, which transforms by
the character $d$ and has linear coefficient one. Set initially
$v_0=Z^s$. At any primitive $p$th root $\zeta$ the value
$v_0(\zeta-1)$ is independent of the root by tame transitivity.
It belongs to $\mathbb Q_p$ (the coefficients are in $\mathbb Z_p$)
and has valuation one, since $Z/X$ is a unit. Rescale $v_0$ by
the $\mathbb Z_p$-unit $-p/v_0(\zeta-1)$, calling the result $v$.
Then $E=v+p$ is a unit times $Q=\Phi_p(1+X)$: divide by the
distinct factors $X-(\zeta-1)$ and compare the residue order $s$.
The invariant series on each letter are $O[[v]]$ and
$\varphi(v)=vE^s$ times a unit. The procyclic generator $\gamma$,
$u=\chi(\gamma)$, $v_p(u-1)=1$, satisfies $\gamma(v)-v\in vE O[[v]]$.
Indeed the difference divided by $v$ vanishes at every primitive
root, so is divisible by $Q$, and the quotient by $E$ is invariant.
Consequently $\gamma(E)/E$ is a unit congruent to one modulo $v$.

Write $N(T)$ here for the product over letters of the
$\Delta$-invariants of the given Wach module. W79H1--W79H3
at $X=0$ show that the tame group acts trivially modulo $X$
on the original module: the regular Taylor action is substitution
on coefficients with fixed crystalline constants. Averaging a
Wach basis as in A81R7 then gives an invariant basis over
$O[[v]]$, with $\gamma$ acting trivially modulo $v$.
The height bound $E^h\varphi N\subset N$ follows from W79H4.

Set $w(k)=k+v_p(k!)$, $A=\lceil2h/s\rceil$ and
$b_W=w(2A-1)$. Then $p^{b_W}N(T_2)\subseteq fN(T_1)$.
Indeed A81R8's Frobenius coefficient argument uses precisely
$\varphi(X)=XQ$, the two Wach heights $h$, and surjectivity on
the residual \'etale modules. All three hold here, so $X^{2h}$
kills the full Wach cokernel. Thus $v^A$ kills its invariant
cokernel by averaging a preimage. On successive $v^j$ quotients
the generator acts by $u^{sj}$. The two polynomials with root
indices $0\le j<A$ and $A\le j<2A$ kill $v^A N(T_2)$
modulo $v^A fN(T_1)$ as in A81R9. Lagrange interpolation on
these $2A$ roots costs $p^{b_W}$: each denominator has valuation
$2A-1+v_p(k!)+v_p((2A-1-k)!)\le b_W$ at index $k$.
Cancellation of $v^A$ gives the assertion. The same estimate
holds on integral extensions by a trivial quotient, whose new
jump is zero. Here $h\ge1$ by the family size conditions.

A bounded target pair has $y\in E^{-h}N(T_2)$, $z\in vN(T_2)$
and $(\gamma-1)y=(\varphi-1)z$. For every integral finite
target cocycle $c$ there exists such a pair representing
$p^{b_W}[c]$ on restriction to $F=K(\mu_p)$: lift $p^{b_W}$
times the fixed generator in the invariant Wach lattice of the
crystalline extension defined by $c$, then take its Frobenius
and generator differences. The kernel identifications of A81R10
and B81L1 apply here by their height-one argument (strict
primitive Taylor lattice, full rank at zero and on residual
\'etale modules); taking $\Delta$-invariants preserves kernels.
Thus these differences have the stated bounds. On the Witt tensor
the lifted generator minus $p^{b_W}$ times the chosen quotient
lift gives the two exact primitive identities of B81L3.
The rational finite-class surjectivity proof there, over $K$,
uses strict filtered exactness and the crystalline dimension
formula. Its scaled integral extension lattice, followed by the
invariant Wach lift just proved, therefore gives a bounded
source pair above $p^k(y,z)$ for some $k$. This auxiliary power
does not enter the final estimate.
% END L81S2

% BEGIN L81S3
\paragraph{L81S3. The divided difference on the kernel pole lattice.}
For the kernel write $N=N(T_0)$, $N_Q=E^{-h}N$ and
$\delta=(\gamma-1)/v$. Both $N,N_Q$ are stable under $\delta$:
the congruences modulo $v$ on $N$ and $\gamma(E)/E$ prove
this directly. Write $U\subset[-h,h]$ for the source's jump set
and $J=((U+\{1,\ldots,h\})\setminus\{0\})\cup\{1,\ldots,h\}$.
Set $c_j=(1-u^{-j})/p\in\mathbb Z_p$ and
$P(B)=\prod_{j\in J}(B-c_j)$. Then
\[
                    \delta P(\delta)N_Q\subseteq N .
\]
The quotient $N_Q/N$ is finite free over $O$ by division by
the monic polynomial $(v+p)^h$, so it is enough to test its
rational Taylor jet at $v=-p$ on every letter. The Taylor
coordinate $1+X=\zeta\exp(Y)$ is equivalent there, since
$E\sim Q$ has a simple zero at $\zeta-1$. The invariant basis
from L81S2 is also a basis before taking invariants, so the
filtered Taylor lattice in these coordinates is exactly that of
W79H2--W79H3. In a fixed filtration-adapted period basis each
line of jump $a$ has the basis $Y^{-a}$ times its period, and the
$h$ pole degrees have exponents $-a-1,\ldots,-a-h$.
On the degree $-j$ term, substitution by $uY$ and multiplication
by $v(Y)^{-1}$ in $\delta$ give diagonal coefficient
$(u^{-j}-1)/(-p)=c_j$ and only higher-degree additional terms.
This triangular operator preserves each period line, where its
diagonal entries are pairwise distinct over $L'$ because $u$
has infinite order. It is diagonalizable on that line.
Strictness puts its $a$ among the jumps of $T_1$, so all these
diagonal entries are roots of $BP(B)$, including the possible
zero. This proves the inclusion without multiplying the constant
denominators from the Taylor change of basis.

Put $k_0=\lfloor h/s\rfloor$, $C=\prod_{j=1}^{k_0}(u^{sj}-1)$.
The graded $v^j$ action gives
$CvN\subseteq(\gamma-1)vN+v^{k_0+1}N$ by the polynomial
argument of B81L4. On $v^{k_0+1}N$ the operator $1-\varphi$
is invertible by geometric series: applying Frobenius contributes
an extra $E^{s(k_0+1)-h}$ with positive exponent, and $(p,v)$
is Frobenius-stable. Both calculations include $k_0=0$.
% END L81S3

% BEGIN L81S4
\paragraph{L81S4. Correcting the divided integral defect.}
On $N_Q\oplus vN$ form the divided defect
\[
           D(y,z)=\delta y-(\varphi-1)z/v\in N_Q .
\]
The inclusion follows from $\varphi(v)/v\sim E^s$. Suppose
$p^k\omega=D(y,z)$ with $\omega\in N_Q$. Write
$Cz=(\gamma-1)x+z_1$, $x\in vN$, $z_1\in v^{k_0+1}N$
by L81S3. Commuting before division by $v$ gives
\[
 Cp^k\omega=\delta\big(Cy-(\varphi-1)x\big)
                      -(\varphi-1)z_1/v .
\]
The first argument lies in $N_Q$ and the last term in $N$.
L81S3 and $\delta N\subset N$ give $Cp^k P(\delta)\omega\in N$.
Torsion freeness of $N_Q/N$ yields $P(\delta)\omega\in N$.
For each $n\in N$, apply the regulator to $vn$ to write
$Cvn=(\gamma-1)x'+z'_1$ with $x'\in vN$ and
$z'_1\in v^{k_0+1}N$. Invert $1-\varphi$ on the last term
and divide the equation by $v$. Thus $Cn\in\mathrm{im}\,D$.
It follows that $CP(\delta)\omega$ lies in this image.
Since $P(B)-P(0)$ is divisible by $B$, its contribution lies in
$\delta N_Q$ and hence in $\mathrm{im}\,D$. Therefore
$p^{d_U}\omega\in\mathrm{im}\,D$, with
\[
             d_U=v_p(CP(0))=w(k_0)+\sum_{j\in J}v_p(j).
\]
Here $v_p(u^{sj}-1)=1+v_p(j)$, whereas division by $p$ in
$c_j$ leaves only $v_p(j)$.

For L81S2's target pair lift $p^{b_W}E^hy$ and $p^{b_W}z/v$
through the invariant Wach map, then divide or multiply back.
The divided defect of these lifts lies in the kernel $N_Q$;
the auxiliary bounded lift over $p^k(y,z)$ makes its $p^k$
multiple a kernel defect by subtraction as in B81L5.
The calculation above corrects it at exactly the allowed cost
$p^{d_U}$. We obtain an invariant bounded source pair mapping
to $p^{b_W+d_U}(y,z)$. Its restriction to $G_F$ has poles at
most $Q^h$ and second coordinate in $XN$, so B81L2 constructs
a crystalline integral extension there. B81L3's Witt identity
shows its class maps to $p^{2b_W+d_U}[c]|_{G_F}$.
Corestriction and division by the unit $[F:K]=p-1$ yield an
integral finite lift of $p^{2b_W+d_U}[c]$ over $K$.

Finally $J\subset[-h,2h]\setminus\{0\}$ gives
$\sum_{j\in J} v_p(j)\le v_p(h!)+v_p((2h)!)\le3h/s$.
Legendre's bound gives $w(m)\le mp/s$, $2A-1\le4h/s+1$,
and $k_0\le h/s$. Consequently
\[
 2b_W+d_U\le(9p/s^2+3/s)h+2p/s
            \le25h/p+3\le e_S ,
\]
using $p/s\le3/2$ and L81S1's ceiling at the last step.
Multiplication by the remaining power proves both
$\mathrm{LiftExponentColumn}$ and $\mathrm{LiftExponentDualPlaces}$
at the stored $e_S$ on every guarded split index. This establishes
the two lifting fields of the split declaration; the inertia and
gluing conditions for its global index inequality are separate.
% END L81S4

\section{The split rows and the global contribution from the columns}

% BEGIN G81S1
\paragraph{G81S1. The two finite ranks on a conjugate pair.}
Use the hub of F78D2 at a twisted split middle-range index $i$ as in
M81B1, with $p=\ell_i$, $\pi\in Z_{\rm good}(i)$. Write
$W=\rho_i(\pi)/\chi_1$ and $W^*=W^\vee(1)$. The declaration
pairs each column $u$ with the conjugate row $w$. Indeed $p\ge5$,
$p$ avoids $q,d,g_0,g_1$, and $-d$ is a nonzero square modulo $p$ by
the split class. Its two roots lift to distinct roots of $X^2+d$
in $\mathbb Q_p$ by Hensel's lemma. The host CM type is distinguished
by the value on $\sqrt{-d}$. At a fixed completion all embeddings
have the same value on this element, since its image is already in
$\mathbb Q_p$. Complex conjugation exchanges the two roots, hence
exchanges the two types of places and preserves their local degrees.
Selecting the type prescribed by $\mathrm{epJ}$ for columns, and
the opposite type for rows, gives a finite disjoint covering and
satisfies both clauses of $\mathrm{hypDecl}$. Thus the witnessed
case of $\mathrm{dcOf}$ applies. Every witness has the same members
by those clauses. S80J1 supplies $\mathrm{ConclusionJa}$, and
the sign selected from it has its own labelled jump identities.
Use those identities for this very sign in the lists of S80J1.

At the column all embeddings have the upward list, and the rank
block has the cloud and the $m$ lows. Its cloud jumps are positive.
At the row every embedding has the downward list; its cloud is
strictly below $-1$, so the cloud of the dual $W^*$ there is
positive. The low weights of $W$ lie outside $[-1,t]$ in either
orientation by the two gaps in S80J1. Thus $0$ is absent from
$W$ at the column and from $W^*$ at the row. Both representations
are crystalline: the sum at the selected point is crystalline
away from the type primes by S80J3, and the rank summand, its
dual, and its Tate twist are crystalline. A fixed rational
vector would give a trivial subrepresentation and hence a jump
at zero by strictness. For either representation the dimension
of its de Rham quotient by $\mathrm{Fil}^0$ is the sum of all
negative jump multiplicities over the embeddings. Only the $m$
lows per embedding can contribute. The crystalline dimension
formula in J80H1 therefore gives
\[
       j_W(u)\le d_um,\qquad
       \dim H^1_f(E_w,W^*)\le d_um,\qquad
       d_u=[E_u:\mathbb Q_p]=[E_w:\mathbb Q_p].
\]
The inequalities hold over every tested coefficient extension.
% END G81S1

% BEGIN G81S2
\paragraph{G81S2. The split row length at the stored exponent.}
Fix the row $w$ conjugate to $u$, any good coefficient field with
integers $O'$, and any of the native surjections onto a
matrix-identified lattice $T$ for $W^*$. Write $k_{O'}$ for
$\operatorname{length}_{O'}(O'/p)$ and $a=X.e_S(i)$.
Every jump of $W^*$ has absolute value at most
$h=\lfloor2b\rfloor+1$ by M81B1 and the dual rule $z\mapsto-1-z$.
At each embedding there are at most $r$ distinct jumps, and none
equals zero. On the unramified field $E_w$, W79H1 and W79H5 give
the de Rham comparison and the good level-one Wach lattice
required by A79L5--A79L8. Apply that invariant theorem to $T$,
which is allowed to be any of the tested lattices. The calculation
of S79P5 uses only $p,h,r$, the absence of zero and the unramified
degree bound. Thus its per-embedding exponent is at most
$2+r(4+2\lfloor\log_p(2h)\rfloor)\le X.c_0(i,u)$, since
the canonical $c_0$ has the same formula at every place. Writing
$V=T[1/p]$, this kills $(V/T)^{G_{E_w}}$ by $p^{X.c_0(i,u)}$.
This module embeds in $p^{-X.c_0(i,u)}T/T$, so its length is at
most $r X.c_0(i,u) k_{O'}$. This also proves the original
$\mathrm{H0ColumnLengthBound}$ for the conjugate row lattice.

Set $P=H^1_f(E_w,T)$, the preimage of the rational finite
subspace. Its torsion is a quotient of $(V/T)^{G_{E_w}}$ by
the invariants sequence for $T\to V\to V/T$. Its free rank
is at most $d_um$ by G81S1: finite generation and
rationalization follow from the integral argument of A81R1.
L81S4 applies the dual-row lifting condition on the whole
family at the prescribed $a$. In particular the image of
finite source cohomology in $P$ contains $p^aP$, since the
target and source of a tested row are permitted subquotients
of the dual sums (the target is a direct summand rationally).
The quotient by that image has length at most
$a d_um k_{O'}+\operatorname{length}_{O'}P_{\rm tors}$,
by the decomposition of a finite module over the DVR into its
free and torsion parts. The cocycle definition of row length
also quotients by integral coboundaries and therefore gives
exactly this cohomological cokernel. Its normalized length is
at most $d_um a+rX.c_0(i,u)$. This proves
$\mathrm{ColengthRowPlaceBound}$ and, assigning this bound to
each conjugate row, $\mathrm{LengthBoundAtPlaces}$ with
$\mathrm{lenRowSplit}$ on the actual split declaration.
% END G81S2

% BEGIN G81S3
\paragraph{G81S3. A regular character corner when roots coincide modulo the prime.}
On this same split index, replace the family matrix at $\pi$ by
the direct sum of its given integral rank and character matrices.
Its rational equivariant isomorphism with the original point
preserves characteristic polynomials, point evaluations and the
Hecke algebra. Elsewhere keep the original matrices. Local tests
on subobjects of rational sums transport by this conjugation.
For row and column lifting the tested lattices in those subobjects
are arbitrary and the standard target rank lattice is still the
original one. The Hecke evaluation ideal and its support are
unchanged by this normalization.

Use D79G1--D79G3 to construct the integral orthogonal projectors
and the rank-$t$ character corner $e_B$ throughout the family.
Since $E_w$ is unramified, inertia surjects onto the full
cyclotomic Galois group over $\mathbb Q_p$: each finite layer
is totally ramified, and compactness gives the inverse limit.
Take $g_0$ with cyclotomic value $1+p$. Its $t$ character
values at $\pi$ are distinct in characteristic zero. Set
$X_0=e_BR(g_0)e_B$. We can perturb this corner integrally away
from $\pi$ to have simple spectrum on its image at every $x$.

To prove this, fix $x\ne\pi$ and extend to an algebraic closure
of the coefficient field. Every absolutely irreducible global
constituent of $R_x$ in a composition series occurs at most once.
Indeed such a finite series, with its constituent isomorphisms,
is defined over a finite coefficient extension: choose bases of
the subspaces and write down their finitely many entries and the
entries of any isomorphism in the algebraic closure. An invariant
subspace over this field is invariant before passing to its
closure by the same matrix equations. Restrict to a fixed place
over $p$ and use the crystalline comparison of S80J3. Exactness
and strictness of the filtered periods give additivity of every
jump multiplicity on the series. A nonzero factor has a jump at
every embedding. Repeating an isomorphic factor would repeat that
jump, contrary to the simple spectrum of Hodge weights in S80J1
(twisting all factors by the same character merely shifts it).

Let $A_x$ be the span algebra of the group matrices over the
algebraic closure. Its semisimple quotient acts on the simple
constituents as the product of their full matrix algebras. This
follows from the structure theorem for finite-dimensional
algebras over an algebraically closed field; every simple
factor is seen on the faithful module, by passage to a composition
series and nilpotence on the resulting filtration. Thus the corner
$e_B A_x e_B$ surjects onto the corresponding product of matrix
corners, by multiplying a lift on both sides by $e_B$. Choose
diagonal matrices with all entries distinct across the nonzero
corner spaces. The filtration of the image of $e_B$ on the
original module has precisely these corner spaces as its successive
quotients, each once. A lift has distinct eigenvalues on that
image. It follows that the discriminant of the characteristic
polynomial there is a nonzero polynomial on the corner.

The Hecke algebra contains a function supported at $x$ with
nonzero value there, by the separating-products proof of G79E1.
After multiplication by this function the integral corners
$e_B R(g)e_B$ still span the entire corner at $x$ over the
algebraic closure. Take finitely many that span. Consequently the
discriminant of $X_0$ plus a linear combination of these supported
corners is a nonzero polynomial in the coefficients. An infinite
integral domain is Zariski dense in the affine space over its
algebraic closure of fractions: induct on the number of variables
using the finite root bound in one variable. Choose coefficients
in the original DVR avoiding zero. Repeat separately for all
$x\ne\pi$, retaining the initial value at $\pi$.

Enlarge the coefficient field once to split all the resulting
polynomials. The eigenvalues are integral because the polynomials
are monic and integral. D79G5 constructs the ordered-root cover
and its free monomial basis over the extended Hecke algebra:
distinctness in the fraction field is enough for injectivity of
evaluation there, and the ordered splitting algebra is torsion-free.
D79G6 produces the integral first-character polynomial $E_1$
of rank one at every point. Write $\delta$ for its product of
root differences evaluated at the distinguished lift of $\pi$.
The ratios at that lift are $(1+p)^{-j}$, $1\le j<t$.
Since $p$ is odd, their differences from one have normalized
valuations $1+v_p(j)$, by the binomial formula for powers of
$1+p$. Hence $v_p(\delta)=v_\delta(i)$ in the native convention
$s_0(i)=1$. All statements here concern the extended coefficient
ring if a splitting field was needed.

For any block write $C_b$ for the global cocycle module constructed
from these matrices as in G79E7. The determinant proof
G79E11--G79E12 gives the evaluation length bounded by the sum
of the block lengths. It only needs rank one off support, the
free ordered-root basis and surjectivity onto the separate blocks,
now supplied above. The off-support assertion of D79G7 persists
after coefficient extension by scalar extension of the common
linear equations for a fixed vector. The extended evaluation
ideal is the old ideal extended to the new DVR: restrictions
to the complement form a finite matrix, and flat extension
preserves its kernel and the image of evaluation at $\pi$.

For the intersection with global coboundaries, repeat the equation
of G79E7 before division by $\delta$. On any block other than the
first character, intertwining gives
$\delta\sum_j b_j(\pi')z(g_j)=0$. For a coboundary
$z(g)=(M(g)-1)v$ the sum equals $-\delta v$, so $\delta^2 v=0$.
In block dimension $d_b$ the coboundaries in $C_b$ consequently
have length at most $2d_b v_\delta(i)k_{O'}$; they are images of
a submodule of the $\delta^2$-torsion of $(L'/O')^{d_b}$.
The first character has zero coboundaries. The total normalized
loss in passing from cocycles to global classes is therefore at
most $2(r+t-1)v_\delta(i)$.
% END G81S3

% BEGIN G81S4
\paragraph{G81S4. The precise remaining column estimate for the split inequality.}
Keep the coefficient ring $O'$ and the global block classes of
G81S3. For the rank block take the kernel of restriction to the
local quotients by $H_f^1$ at all rows. For each character take
the kernel of restriction to all auxiliary inertias in $Q_i$
and to the local quotients by $H_f^1$ at all rows. Denote the
resulting modules by $K_b$. Define $D_{\rm col}$ to be the
image of $\bigoplus_b K_b$ under simultaneous restriction in
each block to the quotients of local cohomology by $H_f^1$ at
all columns. Its normalized length is denoted by $d_{\rm col}$,
allowing infinite length. Write $c$ for the number of columns,
$a=X.e_S(i)$, $v=v_\delta(i)$, and $h_0$ for the source's maximum
normalized invariant length of the nontrivial character ratios
on the columns. A sufficient remaining estimate on these very
modules is
\[
\begin{split}
 d_{\rm col}\le{}&
 \sum_{u\in\mathrm{col}}
        \{a+d_um(a+v)+2v+h_W(i,u,\pi)\}\\
 &+tc h_0+\{4n_V(c+1)-2(r+t-1)\}v .
\end{split}                                                     \tag{*}
\]
This condition involves the common Hecke column modules after the
row and auxiliary restrictions. Neither a global lattice-index
inequality nor an estimate over independently chosen representations
is assumed in the implication below.

To prove the implication, I81G3--I81G4 apply at each row $w$
to any finite list of the rank cuts from G81S3. Their graph
lattice uses the same point complement. The cut presentation by
the columns of $E_1$ in G79E13 uses only their integrality and
intertwining. There is no division by $\delta$ there. G81S2
therefore bounds the normalized loss on this row by
$d_um a+rX.c_0(i,u)$ for its conjugate column $u$.
The class before imposing row finiteness on each character has
auxiliary inertia loss at most $\sum_{P\in Q_i}z_P$ as in G79E8.
At a row apply I81G5 with M81B1--M81B3 to its local cut
presentation. Its normalized loss is at most $(d_w+1)X.c_B(i)$.
This includes primes with $p-1\mid j+1$ on the $j$th block,
since the extra integral $H^2$ term is cyclic as proved in M81B3.
The joint restriction loss is bounded by the sum of the individual
bounds. The kernel after the further column restrictions has
normalized length at most
\[
                    s_\pi+v_p(h_E)+\sum_{0<j<t}u_j .
\]
Indeed the rank classes now lie in the defining Selmer group
as in I81G6. The first character is everywhere unramified;
for the others M81B3 makes the finite local class zero above $p$.
D79G7 and G79E8 ensure inertia vanishes at the other prescribed
places, and the trace to the class group and the cyclotomic
Tate--Shafarevich comparison of G79E6 apply. Here $u_j$ denotes
the native length for the inverse character of exponent $j$.
These are already modules of cohomology classes by G81S3, so a
global coboundary contributes no further loss.

All normalized original terms survive the common coefficient
enlargement. As in I81G2, a finite free basis of $O'$ over the
original $O$ writes the continuous cochains as finitely many
copies before extension; it respects restriction, invariants,
inertia, coboundaries and the crystalline finite parts. Hence
it computes the same base change for the Selmer kernel.
The factor multiplying finite lengths is the relative ramification
index, also the factor multiplying the length of $O/p$.
For positive twists the restriction-of-scalars proof in G79E6
gives exactly their stated $u_j$. The same base change applies
to $z_P$ and to the evaluation ideal by G81S3. Killing and odd
finiteness at the chosen good point supply the required finite
terms and the normalized total auxiliary bound $p^{-2}$.

Add the block lengths using G81S3. Insert (*) and the row bounds,
and replace $d_w+1$ by $d_w+2$ in the nonnegative character
errors. The global coboundary loss uses precisely the allowance
subtracted in (*). For each column the terms involving $a$
become $(2d_um+1)a$ after adding its conjugate row. The remaining
pair terms are $d_um v+2v+h_W(i,u,\pi)+rX.c_0(i,u)$.
The other errors are
\[
    t X.c_B(i)\sum_{w\in\mathrm{row}}(d_w+2)
            +4 n_V(c+1)v+tc h_0 .
\]
By the canonical split identities in DFD8 and N78T2 the sum
of the $(2d_um+1)a$ terms is the exact $\mathrm{sumOver}$ of
$X.\mathrm{len0}(i,\pi)$, and all the remaining errors are
$X.\mathrm{errorBoundHub}(i,\pi)$. This proves the original split
$\mathrm{StarAD}$ inequality at the chosen point under (*).
The rows, the ordinary and dual integral lifting clauses, and
the character condition at each row are supplied above. Bounding
the actual simultaneous image in (*) by the displayed allowance
remains open in this argument.
% END G81S4

\section{The character estimates on a split conjugate pair}

% BEGIN C81S1
\paragraph{C81S1. The two places use a single row allowance.}
Fix the guarded split index and point of G81S1. Keep the good
coefficient extension, $H=\lfloor2b\rfloor$ and
$P=\lfloor\log_p(2H)\rfloor$ of M81B1--M81B3. For the conjugate
pair $u,w$, write $d=[E_u:\mathbb Q_p]=[E_w:\mathbb Q_p]$ and set
\[
 B_0=6+P+7t,\qquad A=t(4+2P),\qquad c(d)=B_0+d n_V A.
\]
At either place the union of the labelled ambient jump lists has
at most $d n_V$ elements. This holds for the original point and
all further coefficient extensions since each of the $d$ labels
has at most $n_V$ jumps. In M81B2 use that cardinality in each
inner sum. The initial degree still costs at most $3+P$, the
$e+1\le t$ terms cost at most $t\{3+d n_V(4+2P)\}$, and the
remaining corrections cost at most $4t+3$. This gives
$\mathrm{cTBY}(Y,H,H,e)\le c(d)$ for every kernel tested at
either place and every $0\le e<t$. M81B2's first cocycle
conclusion thus holds at this smaller integer. M81B3 bounds the
quotient of the $p^{c(d)}$-finite cocycles by the finite
cocycles by $(d+1)c(d)$ in normalized length. I81G5 applies
these assertions to all actual character cuts; the finite graph
lattice, its trivial quotient line and the equivariance on its
kernel work equally at the column. Consequently the combined
normalized loss on this pair is at most $2(d+1)c(d)$ for each
character.

The stored exponent is $c_B=X.c_B(i)=B_0+4n n_V A$, exactly as
in M81B2. Since a column and its conjugate row are distinct,
$2d\le\sum_{v\mid p}[E_v:\mathbb Q_p]=4n$. Hence
$2c(d)\le c_B+B_0$. Moreover $2n_V A\ge B_0$: using
$n_V=r+t\ge2$, $t\ge1$, $P\ge0$, its left side is at least
$16t+8tP\ge6+7t+P$. Thus $dB_0\le4n n_V A$, or
$(d+1)B_0\le c_B$. Combining the inequalities gives
\[
                2(d+1)c(d)\le (d+2)c_B.
\]
This is the original single row's normalized allowance per
character. The estimate includes every split pair, every
$0\le e<t$ and primes with $p-1\mid e+1$.
% END C81S1

% BEGIN C81S2
\paragraph{C81S2. A split criterion involving only the rank block at the columns.}
Let $d_{\rho,\mathrm{col}}$ be the normalized length of the
simultaneous column restriction of the rank-block classes
$K_\rho$ of G81S4, which already satisfy the row conditions.
Retain its notation $a=X.e_S(i)$, $v=v_\delta(i)$,
$c=|\mathrm{col}|$, $h_W(i,u,\pi)$ and $h_0$. Define
\[
 \mathcal S=t\sum_{u\in\mathrm{col}}
       \bigl((d_u+2)c_B-2(d_u+1)c(d_u)\bigr)\ge0 .
\]
It suffices for the split $\mathrm{StarAD}$ inequality at this
point that
\[
\begin{split}
 d_{\rho,\mathrm{col}}\le{}&
 \sum_{u\in\mathrm{col}}
       \{a+d_um(a+v)+2v+h_W(i,u,\pi)\}\\
 &+tc h_0+\{4n_V(c+1)-2(r+t-1)\}v+\mathcal S .
\end{split}
\]
In particular one may drop $\mathcal S$ from this sufficient
upper bound.

Indeed take the global cohomology classes of G81S3 and impose
the auxiliary inertia conditions on the character blocks as in
G81S4. For each of the $t$ characters, impose the local finite
condition at \emph{all} places over $p$ in a single restriction
map. Its image is contained in the product of its projected
images, whose lengths are bounded by C81S1. The total normalized
loss for these restrictions is at most
$t\sum_{u}2(d_u+1)c(d_u)
=t\sum_{w\in\mathrm{row}}(d_w+2)c_B-\mathcal S$.
The character kernels, the auxiliary inertia loss, and the global
coboundaries have exactly the estimates of G81S3--G81S4.
For the rank block first impose all row conditions at cost at
most $\sum_u(d_um a+rX.c_0(i,u))$ by G81S2 and I81G4, and then
impose the columns at cost $d_{\rho,\mathrm{col}}$ by definition.
The remaining rank classes belong to the native Selmer group.
Inserting the displayed criterion makes the two occurrences of
$\mathcal S$ cancel. The pair terms, the global coboundary loss
and the auxiliary contributions now add to the native split
bound exactly as in G81S4, with its original killing and odd
finiteness guards. This accounts for all the character column
images. The displayed estimate for the rank column image remains
to be proved uniformly at the guarded split points.
% END C81S2

\endgroup
